\documentclass[11pt,a4paper]{article}
% ============================================
% GIFT v3.4 — Shared Preamble
% ============================================
% Usage: % ============================================
% GIFT v3.4 — Shared Preamble
% ============================================
% Usage: % ============================================
% GIFT v3.4 — Shared Preamble
% ============================================
% Usage: \input{gift_preamble} after \documentclass

% ENCODING & FONTS
\usepackage[utf8]{inputenc}
\usepackage[T1]{fontenc}
\usepackage{lmodern}

% PAGE LAYOUT (Golden Ratio)
\usepackage[margin=1.618cm, top=2.618cm, bottom=2.618cm]{geometry}

% ESSENTIAL PACKAGES
\usepackage{float}
\usepackage{caption}
\usepackage{subcaption}
\usepackage{setspace}
\usepackage{fancyhdr}
\usepackage{xcolor}
\usepackage{hyperref}
\usepackage{csquotes}
\usepackage{amsmath}
\usepackage{amssymb}
\usepackage{amsthm}
\usepackage{mathtools}
\usepackage{booktabs}
\usepackage{longtable}
\usepackage{array}
\usepackage{multirow}
\usepackage{tikz}
\usepackage{graphicx}
\usepackage{listings}
\usepackage{enumitem}
% Unicode fallbacks (pandoc artifacts that survive markdown→LaTeX)
\DeclareUnicodeCharacter{00B0}{\ensuremath{^\circ}}
\DeclareUnicodeCharacter{221A}{\ensuremath{\sqrt{\,}}}
\DeclareUnicodeCharacter{2713}{\ensuremath{\checkmark}}
\DeclareUnicodeCharacter{2717}{\ensuremath{\times}}
\DeclareUnicodeCharacter{2070}{\ensuremath{^{0}}}
\DeclareUnicodeCharacter{00B9}{\ensuremath{^{1}}}
\DeclareUnicodeCharacter{00B2}{\ensuremath{^{2}}}
\DeclareUnicodeCharacter{00B3}{\ensuremath{^{3}}}
\DeclareUnicodeCharacter{2074}{\ensuremath{^{4}}}
\DeclareUnicodeCharacter{2075}{\ensuremath{^{5}}}
\DeclareUnicodeCharacter{2076}{\ensuremath{^{6}}}
\DeclareUnicodeCharacter{2077}{\ensuremath{^{7}}}
\DeclareUnicodeCharacter{2078}{\ensuremath{^{8}}}
\DeclareUnicodeCharacter{2079}{\ensuremath{^{9}}}
\DeclareUnicodeCharacter{2071}{\ensuremath{^{i}}}
\DeclareUnicodeCharacter{2080}{\ensuremath{_{0}}}
\DeclareUnicodeCharacter{2081}{\ensuremath{_{1}}}
\DeclareUnicodeCharacter{2082}{\ensuremath{_{2}}}
\DeclareUnicodeCharacter{2083}{\ensuremath{_{3}}}
\DeclareUnicodeCharacter{2084}{\ensuremath{_{4}}}
\DeclareUnicodeCharacter{2085}{\ensuremath{_{5}}}
\DeclareUnicodeCharacter{2086}{\ensuremath{_{6}}}
\DeclareUnicodeCharacter{2087}{\ensuremath{_{7}}}
\DeclareUnicodeCharacter{2088}{\ensuremath{_{8}}}
\DeclareUnicodeCharacter{2089}{\ensuremath{_{9}}}
\DeclareUnicodeCharacter{208A}{\ensuremath{_{+}}}
\DeclareUnicodeCharacter{208B}{\ensuremath{_{-}}}
\DeclareUnicodeCharacter{2192}{\ensuremath{\to}}
\DeclareUnicodeCharacter{2502}{|}
\DeclareUnicodeCharacter{25BC}{\ensuremath{\blacktriangledown}}
\DeclareUnicodeCharacter{2500}{\ensuremath{-}}
\DeclareUnicodeCharacter{2514}{\ensuremath{\llcorner}}
\DeclareUnicodeCharacter{251C}{\ensuremath{\vdash}}
\DeclareUnicodeCharacter{2026}{\ldots}
\DeclareUnicodeCharacter{2208}{\ensuremath{\in}}
\DeclareUnicodeCharacter{2229}{\ensuremath{\cap}}
\DeclareUnicodeCharacter{222A}{\ensuremath{\cup}}
\DeclareUnicodeCharacter{2295}{\ensuremath{\oplus}}
\DeclareUnicodeCharacter{2297}{\ensuremath{\otimes}}
\DeclareUnicodeCharacter{2245}{\ensuremath{\cong}}
\DeclareUnicodeCharacter{2248}{\ensuremath{\approx}}
\DeclareUnicodeCharacter{2260}{\ensuremath{\neq}}
\DeclareUnicodeCharacter{2261}{\ensuremath{\equiv}}
\DeclareUnicodeCharacter{2264}{\ensuremath{\leq}}
\DeclareUnicodeCharacter{2265}{\ensuremath{\geq}}
\DeclareUnicodeCharacter{00D7}{\ensuremath{\times}}
\DeclareUnicodeCharacter{00B1}{\ensuremath{\pm}}
\DeclareUnicodeCharacter{2225}{\ensuremath{\|}}
% Greek letters used in body text (pandoc passes them through unchanged)
\DeclareUnicodeCharacter{03B1}{\ensuremath{\alpha}}
\DeclareUnicodeCharacter{03B2}{\ensuremath{\beta}}
\DeclareUnicodeCharacter{03B3}{\ensuremath{\gamma}}
\DeclareUnicodeCharacter{03B4}{\ensuremath{\delta}}
\DeclareUnicodeCharacter{03B5}{\ensuremath{\varepsilon}}
\DeclareUnicodeCharacter{03B6}{\ensuremath{\zeta}}
\DeclareUnicodeCharacter{03B7}{\ensuremath{\eta}}
\DeclareUnicodeCharacter{03B8}{\ensuremath{\theta}}
\DeclareUnicodeCharacter{03B9}{\ensuremath{\iota}}
\DeclareUnicodeCharacter{03BA}{\ensuremath{\kappa}}
\DeclareUnicodeCharacter{03BB}{\ensuremath{\lambda}}
\DeclareUnicodeCharacter{03BC}{\ensuremath{\mu}}
\DeclareUnicodeCharacter{03BD}{\ensuremath{\nu}}
\DeclareUnicodeCharacter{03BE}{\ensuremath{\xi}}
\DeclareUnicodeCharacter{03C0}{\ensuremath{\pi}}
\DeclareUnicodeCharacter{03C1}{\ensuremath{\rho}}
\DeclareUnicodeCharacter{03C3}{\ensuremath{\sigma}}
\DeclareUnicodeCharacter{03C4}{\ensuremath{\tau}}
\DeclareUnicodeCharacter{03C5}{\ensuremath{\upsilon}}
\DeclareUnicodeCharacter{03C6}{\ensuremath{\varphi}}
\DeclareUnicodeCharacter{03C7}{\ensuremath{\chi}}
\DeclareUnicodeCharacter{03C8}{\ensuremath{\psi}}
\DeclareUnicodeCharacter{03C9}{\ensuremath{\omega}}
\DeclareUnicodeCharacter{0394}{\ensuremath{\Delta}}
\DeclareUnicodeCharacter{0398}{\ensuremath{\Theta}}
\DeclareUnicodeCharacter{039B}{\ensuremath{\Lambda}}
\DeclareUnicodeCharacter{03A0}{\ensuremath{\Pi}}
\DeclareUnicodeCharacter{03A3}{\ensuremath{\Sigma}}
\DeclareUnicodeCharacter{03A6}{\ensuremath{\Phi}}
\DeclareUnicodeCharacter{03A8}{\ensuremath{\Psi}}
\DeclareUnicodeCharacter{03A9}{\ensuremath{\Omega}}
\DeclareUnicodeCharacter{2200}{\ensuremath{\forall}}
\DeclareUnicodeCharacter{2203}{\ensuremath{\exists}}
\DeclareUnicodeCharacter{2207}{\ensuremath{\nabla}}
\DeclareUnicodeCharacter{2202}{\ensuremath{\partial}}
\DeclareUnicodeCharacter{2218}{\ensuremath{\circ}}
\DeclareUnicodeCharacter{2032}{\ensuremath{\prime}}
\DeclareUnicodeCharacter{2113}{\ensuremath{\ell}}
\DeclareUnicodeCharacter{210F}{\ensuremath{\hbar}}
\DeclareUnicodeCharacter{27E8}{\ensuremath{\langle}}
\DeclareUnicodeCharacter{27E9}{\ensuremath{\rangle}}
\DeclareUnicodeCharacter{2308}{\ensuremath{\lceil}}
\DeclareUnicodeCharacter{2309}{\ensuremath{\rceil}}
\DeclareUnicodeCharacter{230A}{\ensuremath{\lfloor}}
\DeclareUnicodeCharacter{230B}{\ensuremath{\rfloor}}
\DeclareUnicodeCharacter{222B}{\ensuremath{\int}}
\DeclareUnicodeCharacter{2211}{\ensuremath{\sum}}
\DeclareUnicodeCharacter{220F}{\ensuremath{\prod}}
\DeclareUnicodeCharacter{221E}{\ensuremath{\infty}}
\DeclareUnicodeCharacter{00B5}{\ensuremath{\mu}}
\DeclareUnicodeCharacter{2236}{\ensuremath{\colon}}
\DeclareUnicodeCharacter{2192}{\ensuremath{\to}}
\DeclareUnicodeCharacter{21A6}{\ensuremath{\mapsto}}
\DeclareUnicodeCharacter{21D2}{\ensuremath{\Rightarrow}}
\DeclareUnicodeCharacter{21D4}{\ensuremath{\Leftrightarrow}}
\DeclareUnicodeCharacter{2282}{\ensuremath{\subset}}
\DeclareUnicodeCharacter{2286}{\ensuremath{\subseteq}}
\DeclareUnicodeCharacter{2287}{\ensuremath{\supseteq}}
\DeclareUnicodeCharacter{222B}{\ensuremath{\int}}
\DeclareUnicodeCharacter{2205}{\ensuremath{\emptyset}}
\DeclareUnicodeCharacter{2299}{\ensuremath{\odot}}
\DeclareUnicodeCharacter{2220}{\ensuremath{\angle}}
\DeclareUnicodeCharacter{2135}{\ensuremath{\aleph}}
\DeclareUnicodeCharacter{207B}{\ensuremath{^{-}}}
\DeclareUnicodeCharacter{207A}{\ensuremath{^{+}}}
\DeclareUnicodeCharacter{220E}{\ensuremath{\blacksquare}}
\DeclareUnicodeCharacter{2061}{}
\DeclareUnicodeCharacter{2062}{}
\DeclareUnicodeCharacter{2063}{}
\DeclareUnicodeCharacter{2009}{\,}
\DeclareUnicodeCharacter{200A}{\,}
\DeclareUnicodeCharacter{200B}{}
\DeclareUnicodeCharacter{2212}{\ensuremath{-}}
\DeclareUnicodeCharacter{2194}{\ensuremath{\leftrightarrow}}
\DeclareUnicodeCharacter{2099}{\ensuremath{_{n}}}
\DeclareUnicodeCharacter{2098}{\ensuremath{_{m}}}
\DeclareUnicodeCharacter{2096}{\ensuremath{_{k}}}
\DeclareUnicodeCharacter{2095}{\ensuremath{_{h}}}
\DeclareUnicodeCharacter{2090}{\ensuremath{_{a}}}
\DeclareUnicodeCharacter{2091}{\ensuremath{_{e}}}
\DeclareUnicodeCharacter{2093}{\ensuremath{_{x}}}
\DeclareUnicodeCharacter{1D62}{\ensuremath{_{i}}}
\DeclareUnicodeCharacter{2C7C}{\ensuremath{_{j}}}
\DeclareUnicodeCharacter{2C7B}{\ensuremath{^{e}}}
\DeclareUnicodeCharacter{1D52}{\ensuremath{^{o}}}
\DeclareUnicodeCharacter{1D5}{\ensuremath{^{u}}}
\DeclareUnicodeCharacter{02E1}{\ensuremath{^{l}}}
\DeclareUnicodeCharacter{2295}{\ensuremath{\oplus}}
\DeclareUnicodeCharacter{27FA}{\ensuremath{\Longleftrightarrow}}
\DeclareUnicodeCharacter{27F9}{\ensuremath{\Longrightarrow}}
\DeclareUnicodeCharacter{2299}{\ensuremath{\odot}}
% Provide \tightlist used by pandoc-generated lists (no-op)
\providecommand{\tightlist}{\setlength{\itemsep}{0pt}\setlength{\parskip}{0pt}}
% pandoc longtable column specs: provide \real as passthrough and a
% `none` counter so that calc's \real{x} expansion path works.
\newcounter{none}
\providecommand{\real}[1]{#1}

% LISTINGS
\lstset{
    basicstyle=\small\ttfamily,
    breaklines=true, frame=single,
    keepspaces=true, showstringspaces=false,
    aboveskip=0.8em, belowskip=0.8em
}

% HEADER/FOOTER
\setlength{\headheight}{14pt}
\pagestyle{fancy}
\fancyhf{}
\fancyhead[R]{\thepage}
\renewcommand{\headrulewidth}{0.2pt}

% HYPERREF
\hypersetup{
    colorlinks=true,
    linkcolor=blue, citecolor=blue, urlcolor=blue,
    pdfauthor={Brieuc de La Fourniere}
}

% SPACING
\setstretch{1.2}
\setlength{\parskip}{0.4em}
\setlength{\parindent}{0pt}

% TITLE FORMATTING
\usepackage{titling}
\pretitle{\LARGE\bfseries}
\posttitle{\vspace{-0.4em}}
\preauthor{}
\postauthor{}
\predate{}
\postdate{}
\setlength{\droptitle}{-2.0em}

% CUSTOM COMMANDS — GIFT notation
\newcommand{\E}{\mathrm{E}}
\newcommand{\Gtwo}{\mathrm{G}_2}
\newcommand{\Kseven}{K_7}
\newcommand{\AdS}{\mathrm{AdS}}
\newcommand{\SU}{\mathrm{SU}}
\newcommand{\SO}{\mathrm{SO}}
\newcommand{\U}{\mathrm{U}}
\newcommand{\Sp}{\mathrm{Sp}}
\newcommand{\PSL}{\mathrm{PSL}}
\newcommand{\SM}{\mathrm{SM}}
\newcommand{\Tr}{\mathrm{Tr}}
\newcommand{\rank}{\mathrm{rank}}
\newcommand{\Vol}{\mathrm{Vol}}
\newcommand{\Hol}{\mathrm{Hol}}
\newcommand{\Ric}{\mathrm{Ric}}
\newcommand{\Rm}{\mathrm{Rm}}
\newcommand{\Scal}{\mathrm{Scal}}
\DeclareMathOperator{\diag}{diag}
\DeclareMathOperator{\cond}{cond}
\newcommand{\GIFT}{\textrm{GIFT}}
\newcommand{\CP}{\mathrm{CP}}
\newcommand{\CKM}{\mathrm{CKM}}
\newcommand{\PMNS}{\mathrm{PMNS}}
\newcommand{\EW}{\mathrm{EW}}
\newcommand{\Pl}{\mathrm{Pl}}
\newcommand{\DE}{\mathrm{DE}}
\newcommand{\DM}{\mathrm{DM}}
\newcommand{\KK}{\mathrm{KK}}
\newcommand{\NK}{\mathrm{NK}}
\newcommand{\TCS}{\mathrm{TCS}}
\newcommand{\PINN}{\mathrm{PINN}}
\newcommand{\btwo}{b_2}
\newcommand{\bthree}{b_3}
\newcommand{\typeI}{\textsc{I}}
\newcommand{\typeII}{\textsc{II}}
\newcommand{\typeIII}{\textsc{III}}
\newcommand{\typeIV}{\textsc{IV}}
\newcommand{\proven}{\textsc{Proven}}

% PDF bookmark safety
\pdfstringdefDisableCommands{%
  \def\textsubscript#1{#1}%
  \def\textsuperscript#1{#1}%
  \def\CP{CP}%
  \def\Gtwo{G2}%
  \def\Kseven{K7}%
  \def\GIFT{GIFT}%
  \def\E{E}%
  \def\SM{SM}%
}

% THEOREM ENVIRONMENTS
\newtheorem{theorem}{Theorem}[section]
\newtheorem{proposition}[theorem]{Proposition}
\newtheorem{lemma}[theorem]{Lemma}
\newtheorem{corollary}[theorem]{Corollary}
\theoremstyle{definition}
\newtheorem{definition}[theorem]{Definition}
\theoremstyle{remark}
\newtheorem{remark}[theorem]{Remark}
 after \documentclass

% ENCODING & FONTS
\usepackage[utf8]{inputenc}
\usepackage[T1]{fontenc}
\usepackage{lmodern}

% PAGE LAYOUT (Golden Ratio)
\usepackage[margin=1.618cm, top=2.618cm, bottom=2.618cm]{geometry}

% ESSENTIAL PACKAGES
\usepackage{float}
\usepackage{caption}
\usepackage{subcaption}
\usepackage{setspace}
\usepackage{fancyhdr}
\usepackage{xcolor}
\usepackage{hyperref}
\usepackage{csquotes}
\usepackage{amsmath}
\usepackage{amssymb}
\usepackage{amsthm}
\usepackage{mathtools}
\usepackage{booktabs}
\usepackage{longtable}
\usepackage{array}
\usepackage{multirow}
\usepackage{tikz}
\usepackage{graphicx}
\usepackage{listings}
\usepackage{enumitem}
% Unicode fallbacks (pandoc artifacts that survive markdown→LaTeX)
\DeclareUnicodeCharacter{00B0}{\ensuremath{^\circ}}
\DeclareUnicodeCharacter{221A}{\ensuremath{\sqrt{\,}}}
\DeclareUnicodeCharacter{2713}{\ensuremath{\checkmark}}
\DeclareUnicodeCharacter{2717}{\ensuremath{\times}}
\DeclareUnicodeCharacter{2070}{\ensuremath{^{0}}}
\DeclareUnicodeCharacter{00B9}{\ensuremath{^{1}}}
\DeclareUnicodeCharacter{00B2}{\ensuremath{^{2}}}
\DeclareUnicodeCharacter{00B3}{\ensuremath{^{3}}}
\DeclareUnicodeCharacter{2074}{\ensuremath{^{4}}}
\DeclareUnicodeCharacter{2075}{\ensuremath{^{5}}}
\DeclareUnicodeCharacter{2076}{\ensuremath{^{6}}}
\DeclareUnicodeCharacter{2077}{\ensuremath{^{7}}}
\DeclareUnicodeCharacter{2078}{\ensuremath{^{8}}}
\DeclareUnicodeCharacter{2079}{\ensuremath{^{9}}}
\DeclareUnicodeCharacter{2071}{\ensuremath{^{i}}}
\DeclareUnicodeCharacter{2080}{\ensuremath{_{0}}}
\DeclareUnicodeCharacter{2081}{\ensuremath{_{1}}}
\DeclareUnicodeCharacter{2082}{\ensuremath{_{2}}}
\DeclareUnicodeCharacter{2083}{\ensuremath{_{3}}}
\DeclareUnicodeCharacter{2084}{\ensuremath{_{4}}}
\DeclareUnicodeCharacter{2085}{\ensuremath{_{5}}}
\DeclareUnicodeCharacter{2086}{\ensuremath{_{6}}}
\DeclareUnicodeCharacter{2087}{\ensuremath{_{7}}}
\DeclareUnicodeCharacter{2088}{\ensuremath{_{8}}}
\DeclareUnicodeCharacter{2089}{\ensuremath{_{9}}}
\DeclareUnicodeCharacter{208A}{\ensuremath{_{+}}}
\DeclareUnicodeCharacter{208B}{\ensuremath{_{-}}}
\DeclareUnicodeCharacter{2192}{\ensuremath{\to}}
\DeclareUnicodeCharacter{2502}{|}
\DeclareUnicodeCharacter{25BC}{\ensuremath{\blacktriangledown}}
\DeclareUnicodeCharacter{2500}{\ensuremath{-}}
\DeclareUnicodeCharacter{2514}{\ensuremath{\llcorner}}
\DeclareUnicodeCharacter{251C}{\ensuremath{\vdash}}
\DeclareUnicodeCharacter{2026}{\ldots}
\DeclareUnicodeCharacter{2208}{\ensuremath{\in}}
\DeclareUnicodeCharacter{2229}{\ensuremath{\cap}}
\DeclareUnicodeCharacter{222A}{\ensuremath{\cup}}
\DeclareUnicodeCharacter{2295}{\ensuremath{\oplus}}
\DeclareUnicodeCharacter{2297}{\ensuremath{\otimes}}
\DeclareUnicodeCharacter{2245}{\ensuremath{\cong}}
\DeclareUnicodeCharacter{2248}{\ensuremath{\approx}}
\DeclareUnicodeCharacter{2260}{\ensuremath{\neq}}
\DeclareUnicodeCharacter{2261}{\ensuremath{\equiv}}
\DeclareUnicodeCharacter{2264}{\ensuremath{\leq}}
\DeclareUnicodeCharacter{2265}{\ensuremath{\geq}}
\DeclareUnicodeCharacter{00D7}{\ensuremath{\times}}
\DeclareUnicodeCharacter{00B1}{\ensuremath{\pm}}
\DeclareUnicodeCharacter{2225}{\ensuremath{\|}}
% Greek letters used in body text (pandoc passes them through unchanged)
\DeclareUnicodeCharacter{03B1}{\ensuremath{\alpha}}
\DeclareUnicodeCharacter{03B2}{\ensuremath{\beta}}
\DeclareUnicodeCharacter{03B3}{\ensuremath{\gamma}}
\DeclareUnicodeCharacter{03B4}{\ensuremath{\delta}}
\DeclareUnicodeCharacter{03B5}{\ensuremath{\varepsilon}}
\DeclareUnicodeCharacter{03B6}{\ensuremath{\zeta}}
\DeclareUnicodeCharacter{03B7}{\ensuremath{\eta}}
\DeclareUnicodeCharacter{03B8}{\ensuremath{\theta}}
\DeclareUnicodeCharacter{03B9}{\ensuremath{\iota}}
\DeclareUnicodeCharacter{03BA}{\ensuremath{\kappa}}
\DeclareUnicodeCharacter{03BB}{\ensuremath{\lambda}}
\DeclareUnicodeCharacter{03BC}{\ensuremath{\mu}}
\DeclareUnicodeCharacter{03BD}{\ensuremath{\nu}}
\DeclareUnicodeCharacter{03BE}{\ensuremath{\xi}}
\DeclareUnicodeCharacter{03C0}{\ensuremath{\pi}}
\DeclareUnicodeCharacter{03C1}{\ensuremath{\rho}}
\DeclareUnicodeCharacter{03C3}{\ensuremath{\sigma}}
\DeclareUnicodeCharacter{03C4}{\ensuremath{\tau}}
\DeclareUnicodeCharacter{03C5}{\ensuremath{\upsilon}}
\DeclareUnicodeCharacter{03C6}{\ensuremath{\varphi}}
\DeclareUnicodeCharacter{03C7}{\ensuremath{\chi}}
\DeclareUnicodeCharacter{03C8}{\ensuremath{\psi}}
\DeclareUnicodeCharacter{03C9}{\ensuremath{\omega}}
\DeclareUnicodeCharacter{0394}{\ensuremath{\Delta}}
\DeclareUnicodeCharacter{0398}{\ensuremath{\Theta}}
\DeclareUnicodeCharacter{039B}{\ensuremath{\Lambda}}
\DeclareUnicodeCharacter{03A0}{\ensuremath{\Pi}}
\DeclareUnicodeCharacter{03A3}{\ensuremath{\Sigma}}
\DeclareUnicodeCharacter{03A6}{\ensuremath{\Phi}}
\DeclareUnicodeCharacter{03A8}{\ensuremath{\Psi}}
\DeclareUnicodeCharacter{03A9}{\ensuremath{\Omega}}
\DeclareUnicodeCharacter{2200}{\ensuremath{\forall}}
\DeclareUnicodeCharacter{2203}{\ensuremath{\exists}}
\DeclareUnicodeCharacter{2207}{\ensuremath{\nabla}}
\DeclareUnicodeCharacter{2202}{\ensuremath{\partial}}
\DeclareUnicodeCharacter{2218}{\ensuremath{\circ}}
\DeclareUnicodeCharacter{2032}{\ensuremath{\prime}}
\DeclareUnicodeCharacter{2113}{\ensuremath{\ell}}
\DeclareUnicodeCharacter{210F}{\ensuremath{\hbar}}
\DeclareUnicodeCharacter{27E8}{\ensuremath{\langle}}
\DeclareUnicodeCharacter{27E9}{\ensuremath{\rangle}}
\DeclareUnicodeCharacter{2308}{\ensuremath{\lceil}}
\DeclareUnicodeCharacter{2309}{\ensuremath{\rceil}}
\DeclareUnicodeCharacter{230A}{\ensuremath{\lfloor}}
\DeclareUnicodeCharacter{230B}{\ensuremath{\rfloor}}
\DeclareUnicodeCharacter{222B}{\ensuremath{\int}}
\DeclareUnicodeCharacter{2211}{\ensuremath{\sum}}
\DeclareUnicodeCharacter{220F}{\ensuremath{\prod}}
\DeclareUnicodeCharacter{221E}{\ensuremath{\infty}}
\DeclareUnicodeCharacter{00B5}{\ensuremath{\mu}}
\DeclareUnicodeCharacter{2236}{\ensuremath{\colon}}
\DeclareUnicodeCharacter{2192}{\ensuremath{\to}}
\DeclareUnicodeCharacter{21A6}{\ensuremath{\mapsto}}
\DeclareUnicodeCharacter{21D2}{\ensuremath{\Rightarrow}}
\DeclareUnicodeCharacter{21D4}{\ensuremath{\Leftrightarrow}}
\DeclareUnicodeCharacter{2282}{\ensuremath{\subset}}
\DeclareUnicodeCharacter{2286}{\ensuremath{\subseteq}}
\DeclareUnicodeCharacter{2287}{\ensuremath{\supseteq}}
\DeclareUnicodeCharacter{222B}{\ensuremath{\int}}
\DeclareUnicodeCharacter{2205}{\ensuremath{\emptyset}}
\DeclareUnicodeCharacter{2299}{\ensuremath{\odot}}
\DeclareUnicodeCharacter{2220}{\ensuremath{\angle}}
\DeclareUnicodeCharacter{2135}{\ensuremath{\aleph}}
\DeclareUnicodeCharacter{207B}{\ensuremath{^{-}}}
\DeclareUnicodeCharacter{207A}{\ensuremath{^{+}}}
\DeclareUnicodeCharacter{220E}{\ensuremath{\blacksquare}}
\DeclareUnicodeCharacter{2061}{}
\DeclareUnicodeCharacter{2062}{}
\DeclareUnicodeCharacter{2063}{}
\DeclareUnicodeCharacter{2009}{\,}
\DeclareUnicodeCharacter{200A}{\,}
\DeclareUnicodeCharacter{200B}{}
\DeclareUnicodeCharacter{2212}{\ensuremath{-}}
\DeclareUnicodeCharacter{2194}{\ensuremath{\leftrightarrow}}
\DeclareUnicodeCharacter{2099}{\ensuremath{_{n}}}
\DeclareUnicodeCharacter{2098}{\ensuremath{_{m}}}
\DeclareUnicodeCharacter{2096}{\ensuremath{_{k}}}
\DeclareUnicodeCharacter{2095}{\ensuremath{_{h}}}
\DeclareUnicodeCharacter{2090}{\ensuremath{_{a}}}
\DeclareUnicodeCharacter{2091}{\ensuremath{_{e}}}
\DeclareUnicodeCharacter{2093}{\ensuremath{_{x}}}
\DeclareUnicodeCharacter{1D62}{\ensuremath{_{i}}}
\DeclareUnicodeCharacter{2C7C}{\ensuremath{_{j}}}
\DeclareUnicodeCharacter{2C7B}{\ensuremath{^{e}}}
\DeclareUnicodeCharacter{1D52}{\ensuremath{^{o}}}
\DeclareUnicodeCharacter{1D5}{\ensuremath{^{u}}}
\DeclareUnicodeCharacter{02E1}{\ensuremath{^{l}}}
\DeclareUnicodeCharacter{2295}{\ensuremath{\oplus}}
\DeclareUnicodeCharacter{27FA}{\ensuremath{\Longleftrightarrow}}
\DeclareUnicodeCharacter{27F9}{\ensuremath{\Longrightarrow}}
\DeclareUnicodeCharacter{2299}{\ensuremath{\odot}}
% Provide \tightlist used by pandoc-generated lists (no-op)
\providecommand{\tightlist}{\setlength{\itemsep}{0pt}\setlength{\parskip}{0pt}}
% pandoc longtable column specs: provide \real as passthrough and a
% `none` counter so that calc's \real{x} expansion path works.
\newcounter{none}
\providecommand{\real}[1]{#1}

% LISTINGS
\lstset{
    basicstyle=\small\ttfamily,
    breaklines=true, frame=single,
    keepspaces=true, showstringspaces=false,
    aboveskip=0.8em, belowskip=0.8em
}

% HEADER/FOOTER
\setlength{\headheight}{14pt}
\pagestyle{fancy}
\fancyhf{}
\fancyhead[R]{\thepage}
\renewcommand{\headrulewidth}{0.2pt}

% HYPERREF
\hypersetup{
    colorlinks=true,
    linkcolor=blue, citecolor=blue, urlcolor=blue,
    pdfauthor={Brieuc de La Fourniere}
}

% SPACING
\setstretch{1.2}
\setlength{\parskip}{0.4em}
\setlength{\parindent}{0pt}

% TITLE FORMATTING
\usepackage{titling}
\pretitle{\LARGE\bfseries}
\posttitle{\vspace{-0.4em}}
\preauthor{}
\postauthor{}
\predate{}
\postdate{}
\setlength{\droptitle}{-2.0em}

% CUSTOM COMMANDS — GIFT notation
\newcommand{\E}{\mathrm{E}}
\newcommand{\Gtwo}{\mathrm{G}_2}
\newcommand{\Kseven}{K_7}
\newcommand{\AdS}{\mathrm{AdS}}
\newcommand{\SU}{\mathrm{SU}}
\newcommand{\SO}{\mathrm{SO}}
\newcommand{\U}{\mathrm{U}}
\newcommand{\Sp}{\mathrm{Sp}}
\newcommand{\PSL}{\mathrm{PSL}}
\newcommand{\SM}{\mathrm{SM}}
\newcommand{\Tr}{\mathrm{Tr}}
\newcommand{\rank}{\mathrm{rank}}
\newcommand{\Vol}{\mathrm{Vol}}
\newcommand{\Hol}{\mathrm{Hol}}
\newcommand{\Ric}{\mathrm{Ric}}
\newcommand{\Rm}{\mathrm{Rm}}
\newcommand{\Scal}{\mathrm{Scal}}
\DeclareMathOperator{\diag}{diag}
\DeclareMathOperator{\cond}{cond}
\newcommand{\GIFT}{\textrm{GIFT}}
\newcommand{\CP}{\mathrm{CP}}
\newcommand{\CKM}{\mathrm{CKM}}
\newcommand{\PMNS}{\mathrm{PMNS}}
\newcommand{\EW}{\mathrm{EW}}
\newcommand{\Pl}{\mathrm{Pl}}
\newcommand{\DE}{\mathrm{DE}}
\newcommand{\DM}{\mathrm{DM}}
\newcommand{\KK}{\mathrm{KK}}
\newcommand{\NK}{\mathrm{NK}}
\newcommand{\TCS}{\mathrm{TCS}}
\newcommand{\PINN}{\mathrm{PINN}}
\newcommand{\btwo}{b_2}
\newcommand{\bthree}{b_3}
\newcommand{\typeI}{\textsc{I}}
\newcommand{\typeII}{\textsc{II}}
\newcommand{\typeIII}{\textsc{III}}
\newcommand{\typeIV}{\textsc{IV}}
\newcommand{\proven}{\textsc{Proven}}

% PDF bookmark safety
\pdfstringdefDisableCommands{%
  \def\textsubscript#1{#1}%
  \def\textsuperscript#1{#1}%
  \def\CP{CP}%
  \def\Gtwo{G2}%
  \def\Kseven{K7}%
  \def\GIFT{GIFT}%
  \def\E{E}%
  \def\SM{SM}%
}

% THEOREM ENVIRONMENTS
\newtheorem{theorem}{Theorem}[section]
\newtheorem{proposition}[theorem]{Proposition}
\newtheorem{lemma}[theorem]{Lemma}
\newtheorem{corollary}[theorem]{Corollary}
\theoremstyle{definition}
\newtheorem{definition}[theorem]{Definition}
\theoremstyle{remark}
\newtheorem{remark}[theorem]{Remark}
 after \documentclass

% ENCODING & FONTS
\usepackage[utf8]{inputenc}
\usepackage[T1]{fontenc}
\usepackage{lmodern}

% PAGE LAYOUT (Golden Ratio)
\usepackage[margin=1.618cm, top=2.618cm, bottom=2.618cm]{geometry}

% ESSENTIAL PACKAGES
\usepackage{float}
\usepackage{caption}
\usepackage{subcaption}
\usepackage{setspace}
\usepackage{fancyhdr}
\usepackage{xcolor}
\usepackage{hyperref}
\usepackage{csquotes}
\usepackage{amsmath}
\usepackage{amssymb}
\usepackage{amsthm}
\usepackage{mathtools}
\usepackage{booktabs}
\usepackage{longtable}
\usepackage{array}
\usepackage{multirow}
\usepackage{tikz}
\usepackage{graphicx}
\usepackage{listings}
\usepackage{enumitem}
% Unicode fallbacks (pandoc artifacts that survive markdown→LaTeX)
\DeclareUnicodeCharacter{00B0}{\ensuremath{^\circ}}
\DeclareUnicodeCharacter{221A}{\ensuremath{\sqrt{\,}}}
\DeclareUnicodeCharacter{2713}{\ensuremath{\checkmark}}
\DeclareUnicodeCharacter{2717}{\ensuremath{\times}}
\DeclareUnicodeCharacter{2070}{\ensuremath{^{0}}}
\DeclareUnicodeCharacter{00B9}{\ensuremath{^{1}}}
\DeclareUnicodeCharacter{00B2}{\ensuremath{^{2}}}
\DeclareUnicodeCharacter{00B3}{\ensuremath{^{3}}}
\DeclareUnicodeCharacter{2074}{\ensuremath{^{4}}}
\DeclareUnicodeCharacter{2075}{\ensuremath{^{5}}}
\DeclareUnicodeCharacter{2076}{\ensuremath{^{6}}}
\DeclareUnicodeCharacter{2077}{\ensuremath{^{7}}}
\DeclareUnicodeCharacter{2078}{\ensuremath{^{8}}}
\DeclareUnicodeCharacter{2079}{\ensuremath{^{9}}}
\DeclareUnicodeCharacter{2071}{\ensuremath{^{i}}}
\DeclareUnicodeCharacter{2080}{\ensuremath{_{0}}}
\DeclareUnicodeCharacter{2081}{\ensuremath{_{1}}}
\DeclareUnicodeCharacter{2082}{\ensuremath{_{2}}}
\DeclareUnicodeCharacter{2083}{\ensuremath{_{3}}}
\DeclareUnicodeCharacter{2084}{\ensuremath{_{4}}}
\DeclareUnicodeCharacter{2085}{\ensuremath{_{5}}}
\DeclareUnicodeCharacter{2086}{\ensuremath{_{6}}}
\DeclareUnicodeCharacter{2087}{\ensuremath{_{7}}}
\DeclareUnicodeCharacter{2088}{\ensuremath{_{8}}}
\DeclareUnicodeCharacter{2089}{\ensuremath{_{9}}}
\DeclareUnicodeCharacter{208A}{\ensuremath{_{+}}}
\DeclareUnicodeCharacter{208B}{\ensuremath{_{-}}}
\DeclareUnicodeCharacter{2192}{\ensuremath{\to}}
\DeclareUnicodeCharacter{2502}{|}
\DeclareUnicodeCharacter{25BC}{\ensuremath{\blacktriangledown}}
\DeclareUnicodeCharacter{2500}{\ensuremath{-}}
\DeclareUnicodeCharacter{2514}{\ensuremath{\llcorner}}
\DeclareUnicodeCharacter{251C}{\ensuremath{\vdash}}
\DeclareUnicodeCharacter{2026}{\ldots}
\DeclareUnicodeCharacter{2208}{\ensuremath{\in}}
\DeclareUnicodeCharacter{2229}{\ensuremath{\cap}}
\DeclareUnicodeCharacter{222A}{\ensuremath{\cup}}
\DeclareUnicodeCharacter{2295}{\ensuremath{\oplus}}
\DeclareUnicodeCharacter{2297}{\ensuremath{\otimes}}
\DeclareUnicodeCharacter{2245}{\ensuremath{\cong}}
\DeclareUnicodeCharacter{2248}{\ensuremath{\approx}}
\DeclareUnicodeCharacter{2260}{\ensuremath{\neq}}
\DeclareUnicodeCharacter{2261}{\ensuremath{\equiv}}
\DeclareUnicodeCharacter{2264}{\ensuremath{\leq}}
\DeclareUnicodeCharacter{2265}{\ensuremath{\geq}}
\DeclareUnicodeCharacter{00D7}{\ensuremath{\times}}
\DeclareUnicodeCharacter{00B1}{\ensuremath{\pm}}
\DeclareUnicodeCharacter{2225}{\ensuremath{\|}}
% Greek letters used in body text (pandoc passes them through unchanged)
\DeclareUnicodeCharacter{03B1}{\ensuremath{\alpha}}
\DeclareUnicodeCharacter{03B2}{\ensuremath{\beta}}
\DeclareUnicodeCharacter{03B3}{\ensuremath{\gamma}}
\DeclareUnicodeCharacter{03B4}{\ensuremath{\delta}}
\DeclareUnicodeCharacter{03B5}{\ensuremath{\varepsilon}}
\DeclareUnicodeCharacter{03B6}{\ensuremath{\zeta}}
\DeclareUnicodeCharacter{03B7}{\ensuremath{\eta}}
\DeclareUnicodeCharacter{03B8}{\ensuremath{\theta}}
\DeclareUnicodeCharacter{03B9}{\ensuremath{\iota}}
\DeclareUnicodeCharacter{03BA}{\ensuremath{\kappa}}
\DeclareUnicodeCharacter{03BB}{\ensuremath{\lambda}}
\DeclareUnicodeCharacter{03BC}{\ensuremath{\mu}}
\DeclareUnicodeCharacter{03BD}{\ensuremath{\nu}}
\DeclareUnicodeCharacter{03BE}{\ensuremath{\xi}}
\DeclareUnicodeCharacter{03C0}{\ensuremath{\pi}}
\DeclareUnicodeCharacter{03C1}{\ensuremath{\rho}}
\DeclareUnicodeCharacter{03C3}{\ensuremath{\sigma}}
\DeclareUnicodeCharacter{03C4}{\ensuremath{\tau}}
\DeclareUnicodeCharacter{03C5}{\ensuremath{\upsilon}}
\DeclareUnicodeCharacter{03C6}{\ensuremath{\varphi}}
\DeclareUnicodeCharacter{03C7}{\ensuremath{\chi}}
\DeclareUnicodeCharacter{03C8}{\ensuremath{\psi}}
\DeclareUnicodeCharacter{03C9}{\ensuremath{\omega}}
\DeclareUnicodeCharacter{0394}{\ensuremath{\Delta}}
\DeclareUnicodeCharacter{0398}{\ensuremath{\Theta}}
\DeclareUnicodeCharacter{039B}{\ensuremath{\Lambda}}
\DeclareUnicodeCharacter{03A0}{\ensuremath{\Pi}}
\DeclareUnicodeCharacter{03A3}{\ensuremath{\Sigma}}
\DeclareUnicodeCharacter{03A6}{\ensuremath{\Phi}}
\DeclareUnicodeCharacter{03A8}{\ensuremath{\Psi}}
\DeclareUnicodeCharacter{03A9}{\ensuremath{\Omega}}
\DeclareUnicodeCharacter{2200}{\ensuremath{\forall}}
\DeclareUnicodeCharacter{2203}{\ensuremath{\exists}}
\DeclareUnicodeCharacter{2207}{\ensuremath{\nabla}}
\DeclareUnicodeCharacter{2202}{\ensuremath{\partial}}
\DeclareUnicodeCharacter{2218}{\ensuremath{\circ}}
\DeclareUnicodeCharacter{2032}{\ensuremath{\prime}}
\DeclareUnicodeCharacter{2113}{\ensuremath{\ell}}
\DeclareUnicodeCharacter{210F}{\ensuremath{\hbar}}
\DeclareUnicodeCharacter{27E8}{\ensuremath{\langle}}
\DeclareUnicodeCharacter{27E9}{\ensuremath{\rangle}}
\DeclareUnicodeCharacter{2308}{\ensuremath{\lceil}}
\DeclareUnicodeCharacter{2309}{\ensuremath{\rceil}}
\DeclareUnicodeCharacter{230A}{\ensuremath{\lfloor}}
\DeclareUnicodeCharacter{230B}{\ensuremath{\rfloor}}
\DeclareUnicodeCharacter{222B}{\ensuremath{\int}}
\DeclareUnicodeCharacter{2211}{\ensuremath{\sum}}
\DeclareUnicodeCharacter{220F}{\ensuremath{\prod}}
\DeclareUnicodeCharacter{221E}{\ensuremath{\infty}}
\DeclareUnicodeCharacter{00B5}{\ensuremath{\mu}}
\DeclareUnicodeCharacter{2236}{\ensuremath{\colon}}
\DeclareUnicodeCharacter{2192}{\ensuremath{\to}}
\DeclareUnicodeCharacter{21A6}{\ensuremath{\mapsto}}
\DeclareUnicodeCharacter{21D2}{\ensuremath{\Rightarrow}}
\DeclareUnicodeCharacter{21D4}{\ensuremath{\Leftrightarrow}}
\DeclareUnicodeCharacter{2282}{\ensuremath{\subset}}
\DeclareUnicodeCharacter{2286}{\ensuremath{\subseteq}}
\DeclareUnicodeCharacter{2287}{\ensuremath{\supseteq}}
\DeclareUnicodeCharacter{222B}{\ensuremath{\int}}
\DeclareUnicodeCharacter{2205}{\ensuremath{\emptyset}}
\DeclareUnicodeCharacter{2299}{\ensuremath{\odot}}
\DeclareUnicodeCharacter{2220}{\ensuremath{\angle}}
\DeclareUnicodeCharacter{2135}{\ensuremath{\aleph}}
\DeclareUnicodeCharacter{207B}{\ensuremath{^{-}}}
\DeclareUnicodeCharacter{207A}{\ensuremath{^{+}}}
\DeclareUnicodeCharacter{220E}{\ensuremath{\blacksquare}}
\DeclareUnicodeCharacter{2061}{}
\DeclareUnicodeCharacter{2062}{}
\DeclareUnicodeCharacter{2063}{}
\DeclareUnicodeCharacter{2009}{\,}
\DeclareUnicodeCharacter{200A}{\,}
\DeclareUnicodeCharacter{200B}{}
\DeclareUnicodeCharacter{2212}{\ensuremath{-}}
\DeclareUnicodeCharacter{2194}{\ensuremath{\leftrightarrow}}
\DeclareUnicodeCharacter{2099}{\ensuremath{_{n}}}
\DeclareUnicodeCharacter{2098}{\ensuremath{_{m}}}
\DeclareUnicodeCharacter{2096}{\ensuremath{_{k}}}
\DeclareUnicodeCharacter{2095}{\ensuremath{_{h}}}
\DeclareUnicodeCharacter{2090}{\ensuremath{_{a}}}
\DeclareUnicodeCharacter{2091}{\ensuremath{_{e}}}
\DeclareUnicodeCharacter{2093}{\ensuremath{_{x}}}
\DeclareUnicodeCharacter{1D62}{\ensuremath{_{i}}}
\DeclareUnicodeCharacter{2C7C}{\ensuremath{_{j}}}
\DeclareUnicodeCharacter{2C7B}{\ensuremath{^{e}}}
\DeclareUnicodeCharacter{1D52}{\ensuremath{^{o}}}
\DeclareUnicodeCharacter{1D5}{\ensuremath{^{u}}}
\DeclareUnicodeCharacter{02E1}{\ensuremath{^{l}}}
\DeclareUnicodeCharacter{2295}{\ensuremath{\oplus}}
\DeclareUnicodeCharacter{27FA}{\ensuremath{\Longleftrightarrow}}
\DeclareUnicodeCharacter{27F9}{\ensuremath{\Longrightarrow}}
\DeclareUnicodeCharacter{2299}{\ensuremath{\odot}}
% Provide \tightlist used by pandoc-generated lists (no-op)
\providecommand{\tightlist}{\setlength{\itemsep}{0pt}\setlength{\parskip}{0pt}}
% pandoc longtable column specs: provide \real as passthrough and a
% `none` counter so that calc's \real{x} expansion path works.
\newcounter{none}
\providecommand{\real}[1]{#1}

% LISTINGS
\lstset{
    basicstyle=\small\ttfamily,
    breaklines=true, frame=single,
    keepspaces=true, showstringspaces=false,
    aboveskip=0.8em, belowskip=0.8em
}

% HEADER/FOOTER
\setlength{\headheight}{14pt}
\pagestyle{fancy}
\fancyhf{}
\fancyhead[R]{\thepage}
\renewcommand{\headrulewidth}{0.2pt}

% HYPERREF
\hypersetup{
    colorlinks=true,
    linkcolor=blue, citecolor=blue, urlcolor=blue,
    pdfauthor={Brieuc de La Fourniere}
}

% SPACING
\setstretch{1.2}
\setlength{\parskip}{0.4em}
\setlength{\parindent}{0pt}

% TITLE FORMATTING
\usepackage{titling}
\pretitle{\LARGE\bfseries}
\posttitle{\vspace{-0.4em}}
\preauthor{}
\postauthor{}
\predate{}
\postdate{}
\setlength{\droptitle}{-2.0em}

% CUSTOM COMMANDS — GIFT notation
\newcommand{\E}{\mathrm{E}}
\newcommand{\Gtwo}{\mathrm{G}_2}
\newcommand{\Kseven}{K_7}
\newcommand{\AdS}{\mathrm{AdS}}
\newcommand{\SU}{\mathrm{SU}}
\newcommand{\SO}{\mathrm{SO}}
\newcommand{\U}{\mathrm{U}}
\newcommand{\Sp}{\mathrm{Sp}}
\newcommand{\PSL}{\mathrm{PSL}}
\newcommand{\SM}{\mathrm{SM}}
\newcommand{\Tr}{\mathrm{Tr}}
\newcommand{\rank}{\mathrm{rank}}
\newcommand{\Vol}{\mathrm{Vol}}
\newcommand{\Hol}{\mathrm{Hol}}
\newcommand{\Ric}{\mathrm{Ric}}
\newcommand{\Rm}{\mathrm{Rm}}
\newcommand{\Scal}{\mathrm{Scal}}
\DeclareMathOperator{\diag}{diag}
\DeclareMathOperator{\cond}{cond}
\newcommand{\GIFT}{\textrm{GIFT}}
\newcommand{\CP}{\mathrm{CP}}
\newcommand{\CKM}{\mathrm{CKM}}
\newcommand{\PMNS}{\mathrm{PMNS}}
\newcommand{\EW}{\mathrm{EW}}
\newcommand{\Pl}{\mathrm{Pl}}
\newcommand{\DE}{\mathrm{DE}}
\newcommand{\DM}{\mathrm{DM}}
\newcommand{\KK}{\mathrm{KK}}
\newcommand{\NK}{\mathrm{NK}}
\newcommand{\TCS}{\mathrm{TCS}}
\newcommand{\PINN}{\mathrm{PINN}}
\newcommand{\btwo}{b_2}
\newcommand{\bthree}{b_3}
\newcommand{\typeI}{\textsc{I}}
\newcommand{\typeII}{\textsc{II}}
\newcommand{\typeIII}{\textsc{III}}
\newcommand{\typeIV}{\textsc{IV}}
\newcommand{\proven}{\textsc{Proven}}

% PDF bookmark safety
\pdfstringdefDisableCommands{%
  \def\textsubscript#1{#1}%
  \def\textsuperscript#1{#1}%
  \def\CP{CP}%
  \def\Gtwo{G2}%
  \def\Kseven{K7}%
  \def\GIFT{GIFT}%
  \def\E{E}%
  \def\SM{SM}%
}

% THEOREM ENVIRONMENTS
\newtheorem{theorem}{Theorem}[section]
\newtheorem{proposition}[theorem]{Proposition}
\newtheorem{lemma}[theorem]{Lemma}
\newtheorem{corollary}[theorem]{Corollary}
\theoremstyle{definition}
\newtheorem{definition}[theorem]{Definition}
\theoremstyle{remark}
\newtheorem{remark}[theorem]{Remark}

% pandoc >= 3.2.1 emits \pandocbounded around images; provide a fallback.
% Also clamp all images to \linewidth (natural-size figures overflowed the
% right margin -- the "truncated figures" flagged by the round-2 reviews).
% Kept here rather than in gift_preamble.tex, shared with the 3.4 build.
\providecommand{\pandocbounded}[1]{#1}
\makeatletter
\def\maxwidth{\ifdim\Gin@nat@width>\linewidth\linewidth\else\Gin@nat@width\fi}
\makeatother
\setkeys{Gin}{width=\maxwidth,keepaspectratio}
% md headings carry their own literal numbering ("3.2 What ..."); suppress
% the LaTeX section counters to avoid "0.2 1. Introduction" double numbering.
\setcounter{secnumdepth}{0}
% Loosen justification slightly: long technical identifiers (made breakable
% at build time via \allowbreak) still need extra stretch to avoid overfulls.
\emergencystretch=3em
\fancyhead[L]{The $K_7$ Framework v3.5: Supplement S1}
\hypersetup{pdftitle={The K7 Framework v3.5 Supplement S1: Mathematical Foundations}}

\graphicspath{{../../../figures/}}

\title{%
\LARGE\textbf{Supplement S1: Mathematical Foundations}\\[0.3em]
\large $E_8$ Exceptional Lie Algebra, $G_2$ Holonomy Manifolds, and $K_7$ Construction%
}
\author{Brieuc de La Fourni\`ere \\ \textit{Independent researcher, Beaune, France}}
\date{July 2026 \quad\textbar\quad v3.5}

\begin{document}

\maketitle
\noindent\rule{\textwidth}{0.2pt}

\thispagestyle{fancy}

\section{E₈ Exceptional Lie Algebra, G₂ Holonomy Manifolds, and K₇ Construction}\label{eux2088-exceptional-lie-algebra-gux2082-holonomy-manifolds-and-kux2087-construction}

\emph{Complete mathematical foundations for the K₇ framework, presenting E₈ architecture and K₇ manifold construction.}

\textbf{Lean Verification}: 213 certificate conjuncts, 15 classified axioms (A--F taxonomy; of which 4 external data packages, see main §8.1), 146 .lean files (Lean 4.29.0, zero \texttt{sorry})

\begin{center}\rule{\linewidth}{0.5pt}\end{center}

\section{Abstract}\label{abstract}

This supplement presents the mathematical architecture underlying the K₇ framework. Part I develops the E₈ exceptional Lie algebra with the Exceptional Chain theorem. Part II introduces G₂ holonomy manifolds, including the correct characterization of the g₂ subalgebra as the kernel of the Lie derivative map. Part III discusses K₇ manifold topology: the pair (b₂,b₃)=(21,77) does not appear in the standard TCS catalogue, but the Joyce-Karigiannis (JK) Z₂³ orbifold construction realizes it via T³ × K3 / Z₂³ resolution (§8.4). A four-phase computer-assisted audit (V4 symplectic screen, anti-symplectic obstruction, K3 lattice existence via Mukai \cite{44}, Garbagnati-Sarti \cite{45} and Nikulin \cite{51}, Betti formula) closes the topological count exactly; the analytic torsion-free statement and an explicit polynomial Z₂³ realization are deferred. Part IV establishes the algebraic reference form determining det(g) = 65/32, with Joyce's perturbation theorem providing the existence criterion for a nearby torsion-free metric (subject to its hypothesis ‖T‖ \textless{} ε₀ = 0.1). PINN validation confirms near-G₂ holonomy with V₇ projection reduced by 97\%. All algebraic results are formally verified in Lean 4 (including \texttt{Joyce\allowbreak{}Karigiannis\allowbreak{}Construction.\allowbreak{}lean}, introduced at core v3.4.14 and since migrated to the authors' archive; see §8.4).

\begin{center}\rule{\linewidth}{0.5pt}\end{center}

Part 0: The Octonionic Foundation

\section{0. Why This Framework Exists}\label{why-this-framework-exists}

The K₇ framework emerges from a single algebraic fact:

\textbf{The octonions \ensuremath{\mathbb{O}} are the largest normed division algebra.}

The derivation chain proceeds as follows:

\begin{verbatim}
O (octonions, dim 8)
    │
    ▼
Im(O) = ℝ⁷ (imaginary octonions)
    │
    ▼
G₂ = Aut(O) (automorphism group, dim 14)
    │
    ▼
K₇ with G₂ holonomy (explicit certified metric, [A])
    │
    ▼
Topological invariants (b₂ = 21, b₃ = 77)
    │
    ▼
95 observables (33(I) + 19(II) + 21(III) + 22(IV), 55 Lean-certified)
\end{verbatim}

\textbf{Status}: 95 observables across 4 types (v3.4.29 ledger). 33 Type I algebraic, 19 Type II one-step, 21 Type III multi-step, 22 Type IV structural (incl.~6 metric block eigenvalues + Pinčák 2026 \cite{42}). 55/95 Lean-certified (213 conjuncts, 15 classified axioms of which 4 external data packages, 0 sorry). See §4 of main text and Supplement S3 for the complete dataset. The gauge breaking chain (§5 of main text) is certified in \texttt{TCSGauge\allowbreak{}Breaking.\allowbreak{}lean} and \texttt{Gauge\allowbreak{}Bundle\allowbreak{}Data.\allowbreak{}lean}.

\subsection{0.1 The Division Algebra Chain}\label{the-division-algebra-chain}

\begin{longtable}[]{@{}llll@{}}
\toprule\noalign{}
Algebra & Dim & Physics Role & Stops? \\
\midrule\noalign{}
\endhead
\bottomrule\noalign{}
\endlastfoot
ℝ & 1 & Classical mechanics & No \\
ℂ & 2 & Quantum mechanics & No \\
ℍ & 4 & Spin, Lorentz group & No \\
\textbf{\ensuremath{\mathbb{O}}} & \textbf{8} & \textbf{Exceptional structures} & \textbf{Yes} \\
\end{longtable}

The pattern terminates at \ensuremath{\mathbb{O}}. There is no 16-dimensional normed division algebra. The octonions are \emph{the end of the line}.

\subsection{0.2 G₂ as Octonionic Automorphisms}\label{gux2082-as-octonionic-automorphisms}

\textbf{Definition}: G₂ = \{g ∈ GL(\ensuremath{\mathbb{O}}) : g(xy) = g(x)g(y) for all x,y ∈ \ensuremath{\mathbb{O}}\}

\begin{longtable}[]{@{}lll@{}}
\toprule\noalign{}
Property & Value & Role \\
\midrule\noalign{}
\endhead
\bottomrule\noalign{}
\endlastfoot
dim(G₂) & 14 = C(7,2) − C(7,1) = 21 − 7 & Q\_\allowbreak{}Koide numerator \\
Action & Transitive on S⁶ ⊂ Im(\ensuremath{\mathbb{O}}) & Connects all directions \\
Embedding & G₂ ⊂ SO(7) & Preserves φ₀ \\
\end{longtable}

\subsection{0.3 Why dim(K₇) = 7}\label{why-dimkux2087-7}

The dimension 7 is a consequence of the octonionic structure, not an independent choice: - Im(\ensuremath{\mathbb{O}}) has dimension 7 - G₂ acts naturally on ℝ⁷ - A compact 7-manifold with G₂ holonomy provides the geometric realization

In this sense, K₇ is to G₂ what the circle is to U(1).

\subsection{0.4 The Fano Plane: Combinatorial Structure of Im(\ensuremath{\mathbb{O}})}\label{the-fano-plane-combinatorial-structure-of-imux1d546}

The 7 imaginary octonion units form the \textbf{Fano plane} PG(2,2), the smallest projective plane: - 7 points (imaginary units e₁\ldots e₇) - 7 lines (multiplication triples eᵢ × eⱼ = ±eₖ) - 3 points per line

\textbf{Combinatorial counts}: - Point-line incidences: 7 × 3 = 21 = C(7,2) = b₂ - Automorphism group: PSL(2,7) with \textbar PSL(2,7)\textbar{} = 168

\textbf{Numerical observation}: The following arithmetic identity holds: \[(b_3 + \dim(G_2)) + b_3 = 91 + 77 = 168 = |{\rm PSL}(2,7)| = {\rm rank}(E_8) \times b_2\]

Whether this reflects deeper geometric structure connecting gauge and matter sectors, or is an arithmetic coincidence, remains an open question.

\begin{center}\rule{\linewidth}{0.5pt}\end{center}

Part I: E₈ Exceptional Lie Algebra

\section{1. Root System and Dynkin Diagram}\label{root-system-and-dynkin-diagram}

\subsection{1.1 Basic Data}\label{basic-data}

\begin{longtable}[]{@{}lll@{}}
\toprule\noalign{}
Property & Value & Role \\
\midrule\noalign{}
\endhead
\bottomrule\noalign{}
\endlastfoot
Dimension & dim(E₈) = 248 & Gauge DOF \\
Rank & rank(E₈) = 8 & Cartan subalgebra \\
Number of roots & \textbar Φ(E₈)\textbar{} = 240 & E₈ kissing number \\
Root length & √2 & α\_\allowbreak{}s numerator \\
Coxeter number & h = 30 & Icosahedron edges \\
Dual Coxeter number & h∨ = 30 & McKay correspondence \\
\end{longtable}

\subsection{1.2 Root System Construction}\label{root-system-construction}

E₈ root system in ℝ⁸ has 240 roots:

\textbf{Type I (112 roots)}: Permutations and sign changes of (±1, ±1, 0, 0, 0, 0, 0, 0)

\textbf{Type II (128 roots)}: Half-integer coordinates with even minus signs: \[\frac{1}{2}(\pm 1, \pm 1, \pm 1, \pm 1, \pm 1, \pm 1, \pm 1, \pm 1)\]

\textbf{Verification}: 112 + 128 = 240 roots, all length √2.

\textbf{Lean Status}: E₈ Root System \textbf{12/12 COMPLETE}. All theorems proven: - \texttt{D8\_\allowbreak{}roots\_\allowbreak{}card} = 112, \texttt{Half\allowbreak{}Int\_\allowbreak{}roots\_\allowbreak{}card} = 128 - \texttt{E8\_\allowbreak{}roots\_\allowbreak{}card} = 240, \texttt{E8\_\allowbreak{}roots\_\allowbreak{}decomposition} - \texttt{E8\_\allowbreak{}inner\_\allowbreak{}integral}, \texttt{E8\_\allowbreak{}norm\_\allowbreak{}sq\_\allowbreak{}even}, \texttt{E8\_\allowbreak{}sub\_\allowbreak{}closed} - \texttt{E8\_\allowbreak{}basis\_\allowbreak{}generates}: Every lattice vector is integer combination of simple roots (theorem)

\subsection{1.3 Cartan Matrix}\label{cartan-matrix}

\[A_{E_8} = \begin{pmatrix}
2 & 0 & -1 & 0 & 0 & 0 & 0 & 0 \\
0 & 2 & 0 & -1 & 0 & 0 & 0 & 0 \\
-1 & 0 & 2 & -1 & 0 & 0 & 0 & 0 \\
0 & -1 & -1 & 2 & -1 & 0 & 0 & 0 \\
0 & 0 & 0 & -1 & 2 & -1 & 0 & 0 \\
0 & 0 & 0 & 0 & -1 & 2 & -1 & 0 \\
0 & 0 & 0 & 0 & 0 & -1 & 2 & -1 \\
0 & 0 & 0 & 0 & 0 & 0 & -1 & 2
\end{pmatrix}\]

\textbf{Properties}: det(A) = 1 (unimodular), positive definite.

\begin{center}\rule{\linewidth}{0.5pt}\end{center}

\section{2. Weyl Group}\label{weyl-group}

\subsection{2.1 Order and Factorization}\label{order-and-factorization}

\[|W(E_8)| = 696,729,600 = 2^{14} \times 3^5 \times 5^2 \times 7\]

\subsection{2.2 Topological Factorization Theorem}\label{topological-factorization-theorem}

\textbf{Theorem}: The Weyl group order factorizes into the framework's structural constants:

\[|W(E_8)| = p_2^{\dim(G_2)} \times N_{gen}^{Weyl} \times Weyl^{p_2} \times \dim(K_7)\]

\begin{longtable}[]{@{}llll@{}}
\toprule\noalign{}
Factor & Exponent & Value & Origin \\
\midrule\noalign{}
\endhead
\bottomrule\noalign{}
\endlastfoot
2¹⁴ & dim(G₂) = 14 & 16384 & p₂\^{}(holonomy dim) \\
3⁵ & Weyl = 5 & 243 & N\_\allowbreak{}gen\^{}(Weyl factor) \\
5² & p₂ = 2 & 25 & Weyl\^{}(binary) \\
7¹ & 1 & 7 & dim(K₇) \\
\end{longtable}

\textbf{Status}: \textbf{VERIFIED (Lean 4)}: \texttt{weyl\_\allowbreak{}E8\_\allowbreak{}topological\_\allowbreak{}factorization}

\begin{center}\rule{\linewidth}{0.5pt}\end{center}

\subsection{2.3 Triple Derivation of Weyl = 5}\label{triple-derivation-of-weyl-5}

\textbf{Theorem}: The Weyl factor admits three independent derivations from topological invariants.

\subsection{Derivation 1: G₂ Dimensional Ratio}\label{derivation-1-gux2082-dimensional-ratio}

\[\text{Weyl} = \frac{\dim(G_2) + 1}{N_{gen}} = \frac{14 + 1}{3} = \frac{15}{3} = 5\]

\textbf{Interpretation}: The holonomy dimension plus unity, distributed over generations.

\subsection{Derivation 2: Betti Reduction}\label{derivation-2-betti-reduction}

\[\text{Weyl} = \frac{b_2}{N_{gen}} - p_2 = \frac{21}{3} - 2 = 7 - 2 = 5\]

\textbf{Interpretation}: The per-generation Betti contribution minus binary duality.

\subsection{Derivation 3: Exceptional Difference}\label{derivation-3-exceptional-difference}

\[\text{Weyl} = \dim(G_2) - \text{rank}(E_8) - 1 = 14 - 8 - 1 = 5\]

\textbf{Interpretation}: The gap between holonomy dimension and gauge rank, reduced by unity.

\subsection{Unified Identity}\label{unified-identity}

These three derivations establish the \textbf{Weyl Triple Identity}:

\[\boxed{\frac{\dim(G_2) + 1}{N_{gen}} = \frac{b_2}{N_{gen}} - p_2 = \dim(G_2) - \text{rank}(E_8) - 1 = 5}\]

\textbf{Status}: VERIFIED (algebraic identity from framework constants)

\subsection{Verification}\label{verification}

\begin{longtable}[]{@{}lll@{}}
\toprule\noalign{}
Expression & Computation & Result \\
\midrule\noalign{}
\endhead
\bottomrule\noalign{}
\endlastfoot
(dim(G₂) + 1) / N\_\allowbreak{}gen & (14 + 1) / 3 & 5 \\
b₂/N\_\allowbreak{}gen - p₂ & 21/3 - 2 & 5 \\
dim(G₂) - rank(E₈) - 1 & 14 - 8 - 1 & 5 \\
\end{longtable}

\subsection{Significance}\label{significance}

The triple convergence suggests Weyl = 5 is structurally constrained by the E₈ x E₈/G₂/K₇ geometry. It enters:

\begin{enumerate}
\def\labelenumi{\arabic{enumi}.}
\tightlist
\item
  \textbf{det(g) = 65/32}: Via Weyl x (rank(E₈) + Weyl) / 2\^{}Weyl = 5 x 13 / 32
\item
  \textbf{\textbar W(E₈)\textbar{} factorization}: The factor 5² = Weyl\^{}p₂ in prime decomposition
\item
  \textbf{Cosmological ratio}: sqrt(Weyl) = sqrt(5) appears in dark sector density ratios (see main paper §4 and Supplement S3, §3.4 cosmology tables)
\end{enumerate}

\textbf{Status}: VERIFIED (three independent derivations)

\begin{center}\rule{\linewidth}{0.5pt}\end{center}

\section{3. Exceptional Chain}\label{exceptional-chain}

\subsection{3.1 The Pattern}\label{the-pattern}

A pattern connects exceptional algebra dimensions to primes:

\begin{longtable}[]{@{}lllll@{}}
\toprule\noalign{}
Algebra & n & dim(E\_\allowbreak{}n) & Prime & Index \\
\midrule\noalign{}
\endhead
\bottomrule\noalign{}
\endlastfoot
E₆ & 6 & 78 & 13 & prime(6) \\
E₇ & 7 & 133 & 19 & prime(8) = prime(rank(E₈)) \\
E₈ & 8 & 248 & 31 & prime(11) = prime(D\_\allowbreak{}bulk) \\
\end{longtable}

\subsection{3.2 Exceptional Chain Theorem}\label{exceptional-chain-theorem}

\textbf{Theorem}: For n ∈ \{6, 7, 8\}: \[\dim(E_n) = n \times prime(g(n))\]

where g(6) = 6, g(7) = rank(E₈) = 8, g(8) = D\_\allowbreak{}bulk = 11 (= dim(M₄) + dim(K₇) = 4 + 7, the total bulk dimension of M-theory).

\textbf{Proof} (verified in Lean): - E₆: 6 × 13 = 78 ✓ - E₇: 7 × 19 = 133 ✓ - E₈: 8 × 31 = 248 ✓

\textbf{Status}: \textbf{VERIFIED (Lean 4)}: \texttt{exceptional\_\allowbreak{}chain\_\allowbreak{}certified}

\begin{center}\rule{\linewidth}{0.5pt}\end{center}

\section{4. E₈×E₈ Product Structure}\label{eux2088eux2088-product-structure}

\subsection{4.1 Direct Sum}\label{direct-sum}

\begin{longtable}[]{@{}ll@{}}
\toprule\noalign{}
Property & Value \\
\midrule\noalign{}
\endhead
\bottomrule\noalign{}
\endlastfoot
Dimension & 496 = 248 × 2 \\
Rank & 16 = 8 × 2 \\
Roots & 480 = 240 × 2 \\
\end{longtable}

\subsection{4.2 τ Numerator Connection}\label{ux3c4-numerator-connection}

The hierarchy parameter numerator: \[\tau_{num} = 3472 = 7 \times 496 = \dim(K_7) \times \dim(E_8 \times E_8)\]

\textbf{Status}: \textbf{VERIFIED (Lean 4)}: \texttt{tau\_\allowbreak{}num\_\allowbreak{}E8x\allowbreak{}E8}

\subsection{4.3 Binary Duality Parameter}\label{binary-duality-parameter}

\textbf{Triple geometric origin of p₂ = 2}:

\begin{enumerate}
\def\labelenumi{\arabic{enumi}.}
\tightlist
\item
  \textbf{Local}: p₂ = dim(G₂)/dim(K₇) = 14/7 = 2
\item
  \textbf{Global}: p₂ = dim(E₈×E₈)/dim(E₈) = 496/248 = 2
\item
  \textbf{Root}: √2 in E₈ root normalization
\end{enumerate}

\begin{center}\rule{\linewidth}{0.5pt}\end{center}

\section{5. Exceptional Algebras from Octonions}\label{exceptional-algebras-from-octonions}

The foundational role of octonions is established in Part 0. This section details the exceptional algebraic structures that emerge from \ensuremath{\mathbb{O}}.

\subsection{5.1 Exceptional Jordan Algebra J₃(O)}\label{exceptional-jordan-algebra-jux2083o}

\begin{longtable}[]{@{}ll@{}}
\toprule\noalign{}
Property & Value \\
\midrule\noalign{}
\endhead
\bottomrule\noalign{}
\endlastfoot
dim(J₃(O)) & 27 = 3³ \\
dim(J₃(O)₀) & 26 (traceless) \\
\end{longtable}

\textbf{E-series formula}: The dimension 27 itself emerges from the exceptional chain:

\[\dim(J_3(\mathbb{O})) = \frac{\dim(E_8) - \dim(E_6) - \dim(SU_3)}{6} = \frac{248 - 78 - 8}{6} = \frac{162}{6} = 27\]

This shows the Jordan algebra dimension is derivable from the E-series structure.

\textbf{Status}: \textbf{VERIFIED (Lean 4)}: \texttt{j3o\_\allowbreak{}e\_\allowbreak{}series\_\allowbreak{}certificate}

\subsection{5.2 F₄ Connection}\label{fux2084-connection}

F₄ is the automorphism group of J₃(O): \[\dim(F_4) = 52 = p_2^2 \times \alpha_{sum}^B = 4 \times 13\]

\subsection{5.3 Exceptional Differences}\label{exceptional-differences}

\begin{longtable}[]{@{}lll@{}}
\toprule\noalign{}
Difference & Value & Framework expression \\
\midrule\noalign{}
\endhead
\bottomrule\noalign{}
\endlastfoot
dim(E₈) - dim(J₃(O)) & 221 = 13 × 17 & α\_\allowbreak{}B × λ\_\allowbreak{}H\_\allowbreak{}num \\
dim(F₄) - dim(J₃(O)) & 25 = 5² & Weyl² \\
dim(E₆) - dim(F₄) & 26 & dim(J₃(O)₀) \\
\end{longtable}

\textbf{Status}: \textbf{VERIFIED (Lean 4)}: \texttt{exceptional\_\allowbreak{}differences\_\allowbreak{}certified}

\subsection{5.4 Structural Derivation of τ}\label{structural-derivation-of-ux3c4}

The hierarchy parameter τ admits a purely geometric derivation from framework invariants:

\[\tau = \frac{\dim(E_8 \times E_8) \times b_2}{\dim(J_3(\mathbb{O})) \times H^*} = \frac{496 \times 21}{27 \times 99} = \frac{10416}{2673} = \frac{3472}{891}\]

\textbf{Prime factorization}: - Numerator: 3472 = 2⁴ × 7 × 31 = dim(K₇) × dim(E₈×E₈) - Denominator: 891 = 3⁴ × 11 = N\_\allowbreak{}gen⁴ × D\_\allowbreak{}bulk

\textbf{Alternative form}: τ\_\allowbreak{}num = 7 × 496 = dim(K₇) × dim(E₈×E₈) = 3472

This anchors τ to topological and algebraic invariants, establishing it as a geometric constant rather than a free parameter.

\textbf{Status}: \textbf{VERIFIED (Lean 4)}: \texttt{tau\_\allowbreak{}structural\_\allowbreak{}certificate}

\begin{center}\rule{\linewidth}{0.5pt}\end{center}

Part II: G₂ Holonomy Manifolds

\section{6. Definition and Properties}\label{definition-and-properties}

\subsection{6.1 G₂ as Exceptional Holonomy}\label{gux2082-as-exceptional-holonomy}

\begin{longtable}[]{@{}lll@{}}
\toprule\noalign{}
Property & Value & Role \\
\midrule\noalign{}
\endhead
\bottomrule\noalign{}
\endlastfoot
dim(G₂) & 14 & Q\_\allowbreak{}Koide numerator \\
rank(G₂) & 2 & Lie rank \\
Definition & Aut(O) & Octonion automorphisms \\
\end{longtable}

\textbf{Lean Status}: G₂ Cross Product \textbf{9/11} proven: - \texttt{epsilon\_\allowbreak{}antisymm}, \texttt{epsilon\_\allowbreak{}diag}, \texttt{cross\_\allowbreak{}apply} ✓ - \texttt{G2\_\allowbreak{}cross\_\allowbreak{}bilinear}, \texttt{G2\_\allowbreak{}cross\_\allowbreak{}antisymm}, \texttt{cross\_\allowbreak{}self} ✓ - \texttt{G2\_\allowbreak{}cross\_\allowbreak{}norm} (Lagrange identity ‖u×v‖² = ‖u‖²‖v‖² − ⟨u,v⟩²) ✓ - \texttt{reflect\_\allowbreak{}preserves\_\allowbreak{}lattice} (Weyl reflection) ✓ - Remaining: \texttt{cross\_\allowbreak{}is\_\allowbreak{}octonion\_\allowbreak{}structure} (343-case timeout), \texttt{G2\_\allowbreak{}equiv\_\allowbreak{}characterizations}

\subsection{6.2 G₂ as Kernel of the Lie Derivative}\label{gux2082-as-kernel-of-the-lie-derivative}

The G₂ subalgebra of so(7) admits a precise characterization as the stabilizer of the associative 3-form phi₀. For any antisymmetric matrix A in so(7), the Lie derivative of phi₀ is:

\[L_A(\varphi_0)_{ijk} = A_{ia}\varphi_{ajk} + A_{ja}\varphi_{iak} + A_{ka}\varphi_{ija}\]

The g₂ subalgebra consists of all A for which L\_\allowbreak{}A(phi₀) = 0:

\[\mathfrak{g}_2 = \ker(L) = \{A \in \mathfrak{so}(7) : L_A(\varphi_0) = 0\}\]

This yields the decomposition so(7) = g₂ + V₇, where dim(g₂) = 14 and dim(V₇) = 7. The complement V₇ carries the standard 7-dimensional representation of G₂.

In practice, the kernel is computed via singular value decomposition (SVD) of the linear map L: so(7) --\textgreater{} Lambda³(R⁷). The 14 singular vectors with eigenvalue zero span g₂; the 7 singular vectors with nonzero eigenvalue span V₇.

\textbf{Note}: A heuristic construction based on Fano-plane indices does not produce correct g₂ generators (each such generator is approximately 67\% in g₂ and 33\% in V₇). The kernel-based construction is the correct definition and must be used in all numerical computations involving g₂/V₇ decomposition.

\subsection{6.3 Holonomy Classification (Berger)}\label{holonomy-classification-berger}

\begin{longtable}[]{@{}lll@{}}
\toprule\noalign{}
Dimension & Holonomy & Geometry \\
\midrule\noalign{}
\endhead
\bottomrule\noalign{}
\endlastfoot
\textbf{7} & \textbf{G₂} & \textbf{Exceptional} \\
8 & Spin(7) & Exceptional \\
\end{longtable}

\subsection{6.4 Torsion: Definition and Framework Interpretation}\label{torsion-definition-and-framework-interpretation}

\textbf{Mathematical definition}: Torsion measures failure of G₂ structure to be parallel: \[T = \nabla\phi \neq 0\]

For a G₂ structure φ, the intrinsic torsion decomposes into four irreducible G₂-modules:

\[T \in W_1 \oplus W_7 \oplus W_{14} \oplus W_{27}\]

\begin{longtable}[]{@{}lll@{}}
\toprule\noalign{}
Class & Dimension & Characterization \\
\midrule\noalign{}
\endhead
\bottomrule\noalign{}
\endlastfoot
W₁ & 1 & Scalar: dφ = τ₀ ⋆φ \\
W₇ & 7 & Vector: dφ = 3τ₁ ∧ φ \\
W₁₄ & 14 & Co-closed part of d⋆φ \\
W₂₇ & 27 & Traceless symmetric \\
\end{longtable}

\textbf{Total dimension}: 1 + 7 + 14 + 27 = 49 = 7² = dim(K₇)²

The torsion-free condition requires all four classes to vanish simultaneously, a highly constrained state with 49 conditions.

\textbf{Torsion-free condition}: \[\nabla\phi = 0 \Leftrightarrow d\phi = 0 \text{ and } d*\phi = 0\]

\textbf{Framework interpretation}:

\begin{longtable}[]{@{}
  >{\raggedright\arraybackslash}p{(\columnwidth - 4\tabcolsep) * \real{0.3846}}
  >{\raggedright\arraybackslash}p{(\columnwidth - 4\tabcolsep) * \real{0.3462}}
  >{\raggedright\arraybackslash}p{(\columnwidth - 4\tabcolsep) * \real{0.2692}}@{}}
\toprule\noalign{}
\begin{minipage}[b]{\linewidth}\raggedright
Quantity
\end{minipage} & \begin{minipage}[b]{\linewidth}\raggedright
Meaning
\end{minipage} & \begin{minipage}[b]{\linewidth}\raggedright
Value
\end{minipage} \\
\midrule\noalign{}
\endhead
\bottomrule\noalign{}
\endlastfoot
κ\_\allowbreak{}T = 1/61 & Topological \emph{capacity} for torsion & Fixed by K₇ \\
φ\_\allowbreak{}ref & Algebraic reference form & c × φ₀ \\
T\_\allowbreak{}realized & Actual torsion for global solution & Constrained by Joyce \\
\end{longtable}

\textbf{Key insight}: The 33 dimensionless predictions use only topological invariants (b₂, b₃, dim(G₂)) and are independent of the specific torsion realization. The value κ\_\allowbreak{}T = 1/61 defines the geometric bound on deviations from φ\_\allowbreak{}ref.

\textbf{Physical interactions}: Emerge from the geometry of K₇, with deviations delta(phi) from the reference form bounded by topological constraints. This mechanism is theoretical and its detailed treatment lies beyond the scope of this supplement.

\begin{center}\rule{\linewidth}{0.5pt}\end{center}

\section{7. Topological Invariants}\label{topological-invariants}

\subsection{7.1 Derived Constants}\label{derived-constants}

\begin{longtable}[]{@{}lll@{}}
\toprule\noalign{}
Constant & Formula & Value \\
\midrule\noalign{}
\endhead
\bottomrule\noalign{}
\endlastfoot
det(g) & p₂ + 1/(b₂ + dim(G₂) - N\_\allowbreak{}gen) & 65/32 \\
κ\_\allowbreak{}T & 1/(b₃ - dim(G₂) - p₂) & 1/61 \\
sin²θ\_\allowbreak{}W & b₂/(b₃ + dim(G₂)) & 3/13 \\
\end{longtable}

\subsection{7.2 The 61 Decomposition}\label{the-61-decomposition}

\[\kappa_T^{-1} = 61 = \dim(F_4) + N_{gen}^2 = 52 + 9\]

Alternative: \[61 = \Pi(\alpha^2_B) + 1 = 2 \times 5 \times 6 + 1\]

\textbf{Status}: \textbf{VERIFIED (Lean 4)}: \texttt{kappa\_\allowbreak{}T\_\allowbreak{}inv\_\allowbreak{}decomposition}

\subsection{7.3 Spectral Geometry}\label{spectral-geometry}

The Laplace-Beltrami operator on K₇ admits a discrete spectrum with eigenvalues 0 = λ₀ \textless{} λ₁ ≤ λ₂ ≤ \ldots{} The first non-zero eigenvalue λ₁ (spectral gap) characterizes the geometry's rigidity.

\textbf{Bare topological ratio}: The ratio dim(G₂)/H* provides a topological reference scale:

\[\frac{\dim(G_2)}{H^*} = \frac{14}{b_2 + b_3 + 1} = \frac{14}{99} = 0.1414...\]

This is NOT the spectral gap itself, but a topological bound, see below.

\textbf{Analytical spectral gap}: If the TCS decomposition holds with b₂(M₁) = 11, the first eigenvalue of the scalar Laplacian admits a closed-form expression from the 1D Sturm-Liouville reduction on the TCS neck:

\[\boxed{\lambda_1 = \frac{\pi^2}{L^2 \cdot g_{ss}} = \frac{6\pi^2}{25 \cdot (b_2(M_1) + \mathrm{rank}(E_8))} = \frac{6\pi^2}{475} = 0.12467...}\]

where L = 5 is the effective neck domain length (including ACyl tails), g\_\allowbreak{}ss = (b₂(M₁) + rank(E₈))/(3·rank(G₂)) = 19/6 is the seam metric component (topologically determined by torsion minimization), and b₂(M₁) = 11 is the Picard rank of the first building block in the conjectured TCS realization. Numerical verification (Richardson extrapolation, N=800→1600): λ₁ = 0.12461, matching the formula to \textbf{0.08\%}.

The 169 metric parameters collapse to a single topological integer b₂(M₁), plus two group-theoretic constants (rank(E₈) = 8, rank(G₂) = 2). This provides a closed-form expression for a KK mass gap on a compact G₂ manifold.

\textbf{Comparison with topological ratios:}

\begin{longtable}[]{@{}
  >{\raggedright\arraybackslash}p{(\columnwidth - 6\tabcolsep) * \real{0.2444}}
  >{\raggedright\arraybackslash}p{(\columnwidth - 6\tabcolsep) * \real{0.1556}}
  >{\raggedright\arraybackslash}p{(\columnwidth - 6\tabcolsep) * \real{0.4222}}
  >{\raggedright\arraybackslash}p{(\columnwidth - 6\tabcolsep) * \real{0.1778}}@{}}
\toprule\noalign{}
\begin{minipage}[b]{\linewidth}\raggedright
Expression
\end{minipage} & \begin{minipage}[b]{\linewidth}\raggedright
Value
\end{minipage} & \begin{minipage}[b]{\linewidth}\raggedright
Deviation from λ₁
\end{minipage} & \begin{minipage}[b]{\linewidth}\raggedright
Status
\end{minipage} \\
\midrule\noalign{}
\endhead
\bottomrule\noalign{}
\endlastfoot
6π²/475 (analytical) & 0.12467 & 0.00\% (definition) & \textbf{EXACT} \\
NK Richardson & 0.12461 & 0.05\% & Numerical verification \\
14/99 (bare topological) & 0.14141 & +13.4\% & Topological bound only \\
13/99 (with spinor correction) & 0.13131 & +5.3\% & Approximate (see remark) \\
\end{longtable}

\emph{Remark:} The exact value λ₁ = 6π²/475 involves π (from the Sturm-Liouville eigenvalue structure), not a rational number. The product λ₁×H* = 6π²×99/475 = 12.3364, the confirmed universal invariant (see universality below).

\textbf{Lean status}: \texttt{Spectral.\allowbreak{}Physical\allowbreak{}Spectral\allowbreak{}Gap} (28 theorems, zero axioms). \texttt{Spectral.\allowbreak{}Selberg\allowbreak{}Bridge} connects the spectral gap to the mollified Dirichlet polynomial S\_\allowbreak{}w(T) via the Selberg trace formula.

\textbf{Connection to π²}: The spectral gap formula λ₁ = 6π²/475 provides a direct link between the transcendental number π² and topological integers. The near-identity dim(G₂)/√2 ≈ π² (0.30\%) finds its resolution: the eigenvalue explicitly involves π² through the Sturm-Liouville structure, while dim(G₂) enters through g\_\allowbreak{}ss.

\textbf{Universality}: The confirmed universal invariant is \textbf{λ₁ × H* = 12.3364}, holding for all 66 known G₂ manifolds (including those beyond the CHNP 2015 tabulated range, plus prior TCS literature). The universal law is empirical across all known examples; the analytical mechanism connecting the torsion-free condition to universality remains to be derived. See \texttt{canonical/\allowbreak{}scripts/\allowbreak{}construction\_\allowbreak{}classification.\allowbreak{}py} for the full scan.

\subsection{7.4 Continued Fraction Structure}\label{continued-fraction-structure}

The topological ratio dim(G₂)/H* admits a notable continued fraction representation:

\[\frac{14}{99} = [0; 7, 14] = \cfrac{1}{7 + \cfrac{1}{14}}\]

The only integers appearing are \textbf{7 = dim(K₇)} and \textbf{14 = dim(G₂)}, the two fundamental dimensions of the framework's geometry. Note: this is a property of the topological ratio, not of the spectral gap λ₁ = 6π²/475 (which is irrational).

\subsection{7.5 Pell Equation Structure}\label{pell-equation-structure}

The spectral gap parameters satisfy a Pell equation:

\[\boxed{H^{*2} - 50 \times \dim(G_2)^2 = 1}\]

Explicitly: \[99^2 - 50 \times 14^2 = 9801 - 9800 = 1\]

where \(50 = \dim(K_7)^2 + 1 = 49 + 1\).

\textbf{Fundamental unit}: The Pell equation \(x^2 - 50y^2 = 1\) has fundamental solution \((x_0, y_0) = (99, 14)\), giving:

\[\varepsilon = 7 + \sqrt{50}, \quad \varepsilon^2 = 99 + 14\sqrt{50}\]

\textbf{Continued fraction bridge}: The discriminant \(\sqrt{50}\) has periodic continued fraction \(\sqrt{50} = [7; \overline{14}]\) with period 1, where the partial quotients are exactly dim(K₇) = 7 and dim(G₂) = 14. Combined with the selection principle κ = π²/14 (formalized in \texttt{Spectral.\allowbreak{}Selection\allowbreak{}Principle}), this provides an arithmetic link between the Pell structure and the spectral gap.

\textbf{Status}: TOPOLOGICAL (algebraic identity verified in Lean)

\begin{center}\rule{\linewidth}{0.5pt}\end{center}

Part III: K₇ Manifold Construction

\section{8. Twisted Connected Sum Framework}\label{twisted-connected-sum-framework}

\subsection{8.1 TCS Construction}\label{tcs-construction}

The twisted connected sum (TCS) construction provides the primary method for constructing compact G₂ manifolds from asymptotically cylindrical building blocks.

\textbf{Key insight}: G₂ manifolds can be built by gluing two asymptotically cylindrical (ACyl) G₂ manifolds along their cylindrical ends, with the topology controlled by a twist diffeomorphism φ.

\subsection{8.2 Asymptotically Cylindrical G₂ Manifolds}\label{asymptotically-cylindrical-gux2082-manifolds}

\textbf{Definition}: A complete Riemannian 7-manifold (M, g) with G₂ holonomy is asymptotically cylindrical (ACyl) if there exists a compact subset K ⊂ M such that M ~K is diffeomorphic to (T₀, ∞) × N for some compact 6-manifold N.

\subsection{8.3 Topological Classification}\label{topological-classification}

The K₇ framework constructs an explicit G₂ metric on a compact 7-manifold K₇ with Betti numbers (b₂, b₃) = (21, 77), certified by Newton-Kantorovich theorem ({[}A{]}). The classification of this topological type within known construction methods remains open.

\textbf{Topological status (TCS-specific).} The pair (b₂, b₃) = (21, 77) does not appear among the catalogued TCS-constructed compact G₂ manifolds. Orthogonal TCS is excluded by parity (b₂+b₃=98 is even; CHNP Lemma 6.7). A building block scan (2026-04-14) shows that b₂=21 \textgreater{} 20=max ρ(K3) forces at least one semi-Fano (non-Fano) building block, and all generic K3 lattice embeddings are excluded by the parity theorem. Non-orthogonal TCS and extra-twisted connected sums remain open paths within TCS itself. \textbf{The JK orbifold route, however, is constructive (§8.4) and realizes (21, 77) outside TCS via T³ × K3 / Z₂³ resolution.}

\textbf{CHNP 2015 range.} The CHNP 2015 explicit tabulation covered building blocks with ρ ≤ 9, giving b₂ ≤ 18. Reaching b₂=21 within TCS would require at least one semi-Fano block with ρ ≥ 20 (Shioda-Inose), outside the CHNP tabulated range. The Mori-Mukai classification allows semi-Fano threefolds with ρ up to \textasciitilde20; a TCS construction with such high-ρ blocks has not been explicitly carried out and is superseded for our purposes by the JK route below.

\textbf{Lattice-theoretic analysis.} The K3 lattice Λ\_\allowbreak{}\{K3\} = U³ ⊕ E₈(-1)² admits an orthogonal decomposition N₁ ⊕ N₂ ⊕ ⟨2⟩ with rk(N₁) = 11, rk(N₂) = 10, satisfying the proposed Nikulin embedding conditions (coprime discriminants, signature, discriminant form matching). This is consistent with a TCS construction but does not constitute a proof; the Nikulin embedding is conjectured, and full verification requires an explicit construction of the matching data.

\textbf{Global properties}: - Compact 7-manifold (no boundary) - Simply connected: π₁(K₇) = \{1\}, b₁ = 0 - G₂ holonomy: Hol(g*) = G₂ {[}Joyce, Prop. 11.2.3{]} - Ricci-flat: Ric(g) = 0 - Euler characteristic: χ(K₇) = 0

\textbf{Combinatorial connections}: - b₂ = 21 = C(7,2) = edges in complete graph K₇ - b₃ = 77 = C(7,3) + 2 × b₂ = 35 + 42

\textbf{Status}: The Betti numbers (b₂, b₃) = (21, 77) are certified by the NK metric ({[}A{]}). The spectral analysis ({[}B{]}) independently confirms 21 + 77 near-zero eigenvalues consistent with these Betti numbers. Within the TCS catalogue, the pair (21, 77) does not appear. Outside TCS, the JK Z₂³ orbifold route (§8.4) realizes (21, 77) at the topological/lattice level. An explicit closed-form positive G₂-structure ansatz at the neck level, with determinant constraint, hyperkähler rotation, and five-layer Picard-Lefschetz Wirtinger certificate, is established in {[}D{]}. The Ricci-flat Kähler (Calabi--Yau) metric on the Z₂³-equivariant K3 fibre now has an explicit finitely-parametrised closed-form approximation whose order-3 residual (667 parameters) satisfies an interval-rigorous, assumption-free, machine-checked bound Var(log R) ≤ ε₃ = 1309/10⁷ ≈ 1.309 × 10⁻⁴ \textless{} 10⁻³ on a frozen reproducible 4000-point witness (Krawczyk--Rump-certified K3 points, forward interval arithmetic, no floating-point assumption), formalized in Lean 4 (\texttt{K3Closed\allowbreak{}Form\allowbreak{}Witness}, \texttt{native\_\allowbreak{}decide}, zero \texttt{sorry}, zero added axiom). Promoting this sample bound to a bound over the entire K3 surface remains an open constraint-aware global-positivity (Positivstellensatz/SOS) problem. A complete \emph{smooth analytic} construction (explicit polynomial Z₂³ realization on a Picard-rank-1, η²=8 K3, plus JK 2017 analytic gluing) remains open.

\subsection{8.4 Joyce-Karigiannis Z₂³ Construction, realizes (b₂, b₃) = (21, 77)}\label{joyce-karigiannis-zux2082uxb3-construction-realizes-bux2082-bux2083-21-77}

The Joyce-Karigiannis (JK) framework \cite{43} resolves orbifolds T³ × K3 / G with G ⊂ Aut(T³) × Aut(K3) acting with A₁-type singularities via Eguchi-Hanson gluing, yielding compact torsion-free G₂ manifolds. For G = Z₂³ with a specific mixed-parity action, the resolved 7-manifold N has Betti numbers given by

\[b_2(N) = b_2(T^3 \times K3 / G) + b_0(\text{fixed loci}), \qquad b_3(N) = b_3(T^3 \times K3 / G) + b_1(\text{fixed loci}).\]

A four-phase computer-assisted audit (\texttt{canonical/\allowbreak{}scripts/\allowbreak{}jk\_\allowbreak{}*.\allowbreak{}py}, results in \texttt{canonical/\allowbreak{}results/\allowbreak{}jk\_\allowbreak{}*.\allowbreak{}json}, 2026-05-04) closes the count for the framework's signature:

\textbf{Phase 1: V4 symplectic screen on CI(2,2,2).} The diagonal V4 = ⟨s₁, s₂⟩ ⊂ Z₂³ action on the CI(2,2,2) ⊂ ℂP⁵ K3 fiber has 24 raw fixed points, decomposing into \textbf{12 V4-orbits → 12 T³ components} of the fixed locus. Generators commute on the V4-invariant 9-dimensional quadric subspace.

\textbf{Phase 2, anti-symplectic obstruction.} For any diagonal τ on CI(2,2,2) ⊂ ℂP⁵ and the canonical quadric net R, the determinant ratio satisfies det(τ) / det(R) ≡ 1, so no diagonal P⁵-linear map is anti-symplectic. The full Z₂³ realization must use \emph{intrinsic K3 lattice automorphisms}, not P⁵-linear actions.

\textbf{Phase 2b: K3 lattice abstract existence.} On the K3 lattice Λ = U³ ⊕ E₈(-1)², the Nikulin involution σ₁ = E₈-swap is an order-2 isometry with trace 6 and eigenspaces (14, 8). The coinvariant T\_\allowbreak{}\{σ₁\} ≅ E₈(-2) has discriminant 2⁸. \textbf{Mukai 1988} (V4 = Z₂² ⊂ M₂₃) gives a symplectic K3 action \cite{44}; \textbf{Nikulin's classification} of order-two non-symplectic involutions via 2-elementary lattices (invariants (r, a, δ)) \cite{51} contains the triples (r, a, δ) = (11, 7, 1) and (11, 9, 1), with fixed-locus shapes (g, k) = ((22−r−a)/2, (r−a)/2) = (2, 2) and (1, 1) respectively (see also Artebani-Sarti-Taki, arXiv:0903.3481, for the prime-order survey). η² = 8 is representable in each invariant sublattice.

\textbf{Phase 4: Betti formula.} The fixed locus consists of:

\begin{longtable}[]{@{}lllll@{}}
\toprule\noalign{}
Component & Count & b₀ & b₁ & Source \\
\midrule\noalign{}
\endhead
\bottomrule\noalign{}
\endlastfoot
T³ & 12 & 12 & 36 & V4-fixed K3 points (Phase 1) \\
S¹×Σ₂ & 1 & 1 & 5 & τ : (g,k)=(2,2) \\
S¹×ℂP¹ & 2 & 2 & 2 & τ : k=2 rational curves \\
S¹×T² & 3 & 3 & 9 & s₁τ, s₂τ, s₁s₂τ : (g,k)=(1,1) \\
S¹×ℂP¹ & 3 & 3 & 3 & s₁τ, s₂τ, s₁s₂τ : k=1 rational \\
\textbf{Total} & \textbf{21} & \textbf{21} & \textbf{55} & -- \\
\end{longtable}

Combined with the quotient cohomology b₂(T³ × K3 / Z₂³) = 0 and b₃(T³ × K3 / Z₂³) = 22, the JK formula gives

\[b_2(N) = 0 + 21 = 21, \qquad b_3(N) = 22 + 55 = 77, \qquad \chi(N) = 0,\]

matching the framework's topological signature exactly.

\textbf{Lean formalization} (introduced at core v3.4.14): the four-phase audit is encoded in \texttt{Joyce\allowbreak{}Karigiannis\allowbreak{}Construction.\allowbreak{}lean} with master theorem \texttt{jk\_\allowbreak{}z23\_\allowbreak{}construction\_\allowbreak{}realizes\_\allowbreak{}gift\_\allowbreak{}betti} proving \texttt{phase4.\allowbreak{}b2N\ =\ GIFT.\allowbreak{}Core.\allowbreak{}b2\ ∧\ phase4.\allowbreak{}b3N\ =\ GIFT.\allowbreak{}Core.\allowbreak{}b3} by \texttt{native\_\allowbreak{}decide}. The module introduces 0 axioms and 0 sorry; literature dependencies (Mukai; Nikulin's involution classification) are encoded as Bool flags in a \texttt{JKPhase2b\allowbreak{}Lattice\allowbreak{}Screen} structure with explicit citation comments. The module was migrated out of the public core release during a later repository audit and is preserved in the authors' archive (available on request); the phase numbering 1, 2, 2b, 4 follows the Lean encoding (2b refines 2; no phase 3).

\textbf{Honest scope.} This route certifies the \emph{topological/lattice gate} only: - ✓ Betti pair (b₂, b₃) = (21, 77) closed by exact integer arithmetic - ✓ K3 surface with required Z₂³ action exists abstractly (Mukai/G-S) - ⏳ Explicit polynomial coordinate model on a Picard-rank-1, η² = 8 K3, deferred (existence guaranteed but moduli construction not done here) - ⏳ Smooth analytic torsion-free G₂ certificate, deferred to JK 2017 gluing theorem

The full four-phase topological route (lattice screens, explicit (r, a, δ) bookkeeping, and Betti formula closure) follows the Joyce-Karigiannis Z₂³ framework \cite{43}. An explicit closed-form positive G₂-structure ansatz at the neck level, determinant-preserving radial family, hyperkähler rotation, base-coframe absorption, and five-layer Picard-Lefschetz Wirtinger certificate, is established in {[}D{]}.

\textbf{Comparison to TCS within the K₇ framework.} The NK-certified metric ({[}A{]}) and the Donaldson analytic note {[}D{]} operate at the \emph{neck level} of a putative compact extension and do not depend on whether the global construction is TCS or JK. The JK route provides a constructive existence proof for the topological pair; the NK metric and {[}D{]} provide complementary analytic evidence at the neck: the former via Newton-Kantorovich certificate, the latter via explicit closed-form positive G₂-structure ansatz with Wirtinger certificate; together they argue that the framework's signature (b₂, b₃) = (21, 77) is realizable, with the JK route being the leading candidate for the global compact construction.

\textbf{Compatibility with intermediate-curvature splitting (Chen-Hong 2026).} Recent work of Chen and Hong \cite{48} establishes a sharp dichotomy for noncompact manifolds with nonnegative m-intermediate curvature admitting a smooth proper degree-nonzero map to M\^{}\{n−m\} × T\^{}\{m−1\} × ℝ. For 3 ≤ n ≤ 5 with 1 ≤ m ≤ n−1, or for 6 ≤ n ≤ 7 with m ∈ \{1, n−2, n−1\}, the manifold is forced to split isometrically as the Riemannian product E × T\^{}\{m−1\} × ℝ. By contrast, for 6 ≤ n ≤ 7 and 2 ≤ m ≤ n−3, they construct explicit metrics on S\^{}\{n−m\} × T\^{}\{m−1\} × ℝ with uniformly positive m-intermediate curvature, showing the algebraic gate m² − mn + m + n \textgreater{} 0 is sharp. The framework's coassociative neck has topology K3 × T² × ℝ (n = 7, m = 3, giving m² − mn + m + n = −2 \textless{} 0), placing it within the Chen-Hong ``non-rigid'' window: metrics on this topology with nonnegative 3-intermediate curvature need not split. The k ≥ 1 Chebyshev corrections that lift the holonomy from SU(2) × U(1) to full G₂ are therefore consistent with the Chen-Hong existence regime, rather than obstructed by a forced splitting.

\begin{center}\rule{\linewidth}{0.5pt}\end{center}

\section{9. Cohomological Structure}\label{cohomological-structure}

\subsection{9.1 Mayer-Vietoris Analysis}\label{mayer-vietoris-analysis}

The Mayer-Vietoris sequence provides the primary tool for computing cohomology:

\[\cdots \to H^{k-1}(N) \xrightarrow{\delta} H^k(K_7) \xrightarrow{i^*} H^k(M_1) \oplus H^k(M_2) \xrightarrow{j^*} H^k(N) \to \cdots\]

\subsection{9.2 Betti Number Derivation}\label{betti-number-derivation}

\textbf{Result for b₂}: b₂(K₇) = 21 is certified by the NK metric. The spectral analysis ({[}B{]}) independently confirms 21 near-zero eigenvalues of Δ₂ with gap ratio 14,635. A TCS realization would require building blocks with Picard ranks summing to 21, but the specific identification of those blocks is an open problem.

\textbf{Result for b₃}: b₃(K₇) = 77 is certified by the NK metric. The harmonic decomposition b₃ = 35 + 42 = (1+7+27) + 2×21 is confirmed by the certified metric with spectral gap 10522× between zero and non-zero modes. A splitting b₃ = b₃(M₁) + b₃(M₂) via Mayer-Vietoris is conditional on the building block identification, which remains open.

\textbf{Status}: VERIFIED (NK metric, {[}A{]}); candidate TCS decomposition is conjectural

\subsection{9.3 Complete Betti Spectrum and Poincaré Duality}\label{complete-betti-spectrum-and-poincaruxe9-duality}

For a compact G₂-holonomy 7-manifold K₇, Poincaré duality gives b\_\allowbreak{}k = b\_\allowbreak{}\{7-k\}:

\begin{longtable}[]{@{}
  >{\raggedright\arraybackslash}p{(\columnwidth - 4\tabcolsep) * \real{0.1250}}
  >{\raggedright\arraybackslash}p{(\columnwidth - 4\tabcolsep) * \real{0.3750}}
  >{\raggedright\arraybackslash}p{(\columnwidth - 4\tabcolsep) * \real{0.5000}}@{}}
\toprule\noalign{}
\begin{minipage}[b]{\linewidth}\raggedright
k
\end{minipage} & \begin{minipage}[b]{\linewidth}\raggedright
b\_\allowbreak{}k(K₇)
\end{minipage} & \begin{minipage}[b]{\linewidth}\raggedright
Derivation
\end{minipage} \\
\midrule\noalign{}
\endhead
\bottomrule\noalign{}
\endlastfoot
0 & 1 & Connected \\
1 & 0 & Simply connected (G₂ holonomy) \\
2 & 21 & NK-certified; spectrally confirmed ({[}B{]}); building block decomposition open \\
3 & 77 & NK-certified; harmonic decomposition 35+42; building block decomposition open \\
4 & 77 & Poincaré duality: b₄ = b₃ \\
5 & 21 & Poincaré duality: b₅ = b₂ \\
6 & 0 & Poincaré duality: b₆ = b₁ \\
7 & 1 & Poincaré duality: b₇ = b₀ \\
\end{longtable}

\textbf{Euler characteristic}: For any compact oriented odd-dimensional manifold, χ = 0: \[\chi(K_7) = \sum_{k=0}^{7} (-1)^k b_k = 1 - 0 + 21 - 77 + 77 - 21 + 0 - 1 = 0\]

\textbf{Status}: \textbf{VERIFIED (Lean 4)}: \texttt{euler\_\allowbreak{}char\_\allowbreak{}K7\_\allowbreak{}is\_\allowbreak{}zero}, \texttt{poincare\_\allowbreak{}duality\_\allowbreak{}K7}

\textbf{Effective cohomological dimension}: \[H^* = b_2 + b_3 + 1 = 21 + 77 + 1 = 99\]

\subsection{9.4 The Structural Constant 42}\label{the-structural-constant-42}

The number 42 appears throughout the framework as a derived topological invariant:

\[42 = 2 \times 3 \times 7 = p_2 \times N_{gen} \times \dim(K_7)\]

\textbf{Multiple derivations}:

\begin{longtable}[]{@{}lll@{}}
\toprule\noalign{}
Formula & Value & Interpretation \\
\midrule\noalign{}
\endhead
\bottomrule\noalign{}
\endlastfoot
p₂ × N\_\allowbreak{}gen × dim(K₇) & 2 × 3 × 7 = 42 & Binary × generations × fiber \\
2 × b₂ & 2 × 21 = 42 & Twice the gauge moduli \\
b₃ - C(7,3) & 77 - 35 = 42 & Global vs local 3-forms \\
\end{longtable}

\textbf{Connection to b₃ decomposition}: \[b_3 = 77 = C(7,3) + 42 = 35 + 2 \times b_2\]

The 35 local modes correspond to Λ³(ℝ⁷) fiber forms; the 42 global modes arise from the TCS structure.

\textbf{Status}: \textbf{VERIFIED (Lean 4)}: \texttt{structural\_\allowbreak{}42\_\allowbreak{}gift\_\allowbreak{}form}, \texttt{structural\_\allowbreak{}42\_\allowbreak{}from\_\allowbreak{}b2}

\subsection{9.5 Third Betti Number Decomposition}\label{third-betti-number-decomposition}

The b₃ = 77 harmonic 3-forms decompose as:

\[H^3(K_7) = H^3_{\text{local}} \oplus H^3_{\text{global}}\]

\begin{longtable}[]{@{}lll@{}}
\toprule\noalign{}
Component & Dimension & Origin \\
\midrule\noalign{}
\endhead
\bottomrule\noalign{}
\endlastfoot
H³\_\allowbreak{}local & 35 = C(7,3) & Λ³(ℝ⁷) fiber forms \\
H³\_\allowbreak{}global & 42 = 2 × 21 & TCS global modes \\
\end{longtable}

\textbf{Verification}: 35 + 42 = 77

\textbf{Status}: TOPOLOGICAL

\begin{center}\rule{\linewidth}{0.5pt}\end{center}

Part IV: Metric Structure and Verification

\section{10. Structural Metric Invariants}\label{structural-metric-invariants}

\subsection{10.1 Structural Metric Invariants and Normalizations}\label{structural-metric-invariants-and-normalizations}

The K₇ framework explores the hypothesis that metric invariants derive from fixed mathematical structure. The topological constraints serve as inputs; the specific geometry is then determined.

\begin{longtable}[]{@{}
  >{\raggedright\arraybackslash}p{(\columnwidth - 6\tabcolsep) * \real{0.3143}}
  >{\raggedright\arraybackslash}p{(\columnwidth - 6\tabcolsep) * \real{0.2571}}
  >{\raggedright\arraybackslash}p{(\columnwidth - 6\tabcolsep) * \real{0.2000}}
  >{\raggedright\arraybackslash}p{(\columnwidth - 6\tabcolsep) * \real{0.2286}}@{}}
\toprule\noalign{}
\begin{minipage}[b]{\linewidth}\raggedright
Invariant
\end{minipage} & \begin{minipage}[b]{\linewidth}\raggedright
Formula
\end{minipage} & \begin{minipage}[b]{\linewidth}\raggedright
Value
\end{minipage} & \begin{minipage}[b]{\linewidth}\raggedright
Status
\end{minipage} \\
\midrule\noalign{}
\endhead
\bottomrule\noalign{}
\endlastfoot
κ\_\allowbreak{}T & 1/(b₃ - dim(G₂) - p₂) & 1/61 & TOPOLOGICAL \\
det(g) & (Weyl × (rank(E₈) + Weyl))/2⁵ & 65/32 & METRIC NORMALIZATION (see §10.3) \\
\end{longtable}

\subsection{10.2 Torsion Magnitude κ\_\allowbreak{}T = 1/61}\label{torsion-magnitude-ux3ba_t-161}

\textbf{Derivation}: \[\kappa_T = \frac{1}{b_3 - \dim(G_2) - p_2} = \frac{1}{77 - 14 - 2} = \frac{1}{61}\]

\textbf{Interpretation}: - 61 = effective matter degrees of freedom - b₃ = 77 total fermion modes - dim(G₂) = 14 gauge symmetry constraints - p₂ = 2 binary duality factor

\textbf{Status}: TOPOLOGICAL

\subsection{10.3 Metric Determinant det(g) = 65/32}\label{metric-determinant-detg-6532}

The metric determinant admits three independent derivations from topological invariants, providing strong evidence for its structural necessity.

\textbf{Path 1} (Weyl formula): \[\det(g) = \frac{\text{Weyl} \times (\text{rank}(E_8) + \text{Weyl})}{2^{\text{Weyl}}} = \frac{5 \times 13}{32} = \frac{65}{32}\]

\textbf{Path 2} (Cohomological): \[\det(g) = p_2 + \frac{1}{b_2 + \dim(G_2) - N_{\text{gen}}} = 2 + \frac{1}{21+14-3} = 2 + \frac{1}{32} = \frac{65}{32}\]

\textbf{Path 3} (H* formula): \[\det(g) = \frac{H^* - b_2 - 13}{32} = \frac{99 - 21 - 13}{32} = \frac{65}{32}\]

The convergence of three independent algebraic paths to the same rational value is suggestive but does not constitute a derivation from topology. The actual chain is: the metric optimization was constrained to det(g) = 65/32 (a normalization target chosen to be compatible with the known structural parameters g\_\allowbreak{}ss = 19/6, g\_\allowbreak{}T² = 7/6), and the above formulas were identified post-hoc. Any rational number with small numerator/denominator can be expressed as combinations of a few integers.

\textbf{Numerical value}: 65/32 = 2.03125 (exact rational)

\textbf{Status}: STRUCTURAL (metric normalization, algebraically exact in our metric, with three suggestive integer formulas, but not derived from topology). The 6 observables using det\_\allowbreak{}num or det\_\allowbreak{}den (six Type I observables) depend on this normalization choice.

\begin{center}\rule{\linewidth}{0.5pt}\end{center}

\section{11. Formal Certification}\label{formal-certification}

\subsection{11.1 The Algebraic Reference Form}\label{the-algebraic-reference-form}

The algebraic reference form in a local G₂-adapted orthonormal coframe:

\[\varphi_{\text{ref}} = c \cdot \varphi_0, \quad c = \left(\frac{65}{32}\right)^{1/14}\] \[g_{\text{ref}} = c^2 \cdot I_7 = \left(\frac{65}{32}\right)^{1/7} \cdot I_7\]

\textbf{Important clarification}: This representation holds in a local orthonormal frame. The manifold K₇ is curved and compact; ``I₇'' reflects the frame choice, not global flatness. The reference form φ\_\allowbreak{}ref yields det(g) = 65/32. The NK-certified model provides computational/conditional evidence for a nearby torsion-free G₂ structure within the certified analytic setup of Paper A; this does not by itself close the independent smooth compact construction problem for K₇, whose full building-block / gluing realization remains open.

\begin{longtable}[]{@{}
  >{\raggedright\arraybackslash}p{(\columnwidth - 4\tabcolsep) * \real{0.4000}}
  >{\raggedright\arraybackslash}p{(\columnwidth - 4\tabcolsep) * \real{0.2800}}
  >{\raggedright\arraybackslash}p{(\columnwidth - 4\tabcolsep) * \real{0.3200}}@{}}
\toprule\noalign{}
\begin{minipage}[b]{\linewidth}\raggedright
Property
\end{minipage} & \begin{minipage}[b]{\linewidth}\raggedright
Value
\end{minipage} & \begin{minipage}[b]{\linewidth}\raggedright
Status
\end{minipage} \\
\midrule\noalign{}
\endhead
\bottomrule\noalign{}
\endlastfoot
det(g) & 65/32 & exact by imposed normalization, not fitted to experiment (§10.3) \\
φ\_\allowbreak{}ref components & 7/35 & 20\% sparsity \\
Joyce threshold & ‖T‖ \textless{} ε₀ = 0.1 & Satisfied (224× margin) \\
\end{longtable}

\subsection{11.2 Joyce Existence Theorem and Global Solutions}\label{joyce-existence-theorem-and-global-solutions}

\textbf{Important clarification}: The reference form φ\_\allowbreak{}ref = c·φ₀ is the canonical G₂ structure in a local orthonormal coframe, not a globally constant form on K₇. On a compact G₂ manifold, the coframe 1-forms \{eⁱ\} satisfy deⁱ ≠ 0 in general, so ``constant components'' does not imply dφ = 0 globally.

\textbf{Actual solution structure}: The topology and geometry of K₇ impose a deformation: \[\varphi = \varphi_{\text{ref}} + \delta\varphi\]

The torsion-free condition (dφ = 0, d*φ = 0) is a \textbf{global constraint}. Joyce's perturbation theorem provides the existence criterion for a nearby torsion-free G₂ metric, conditional on the initial torsion satisfying ‖T‖ \textless{} ε₀ = 0.1. PINN validation (N=1000) confirms ‖T‖\_\allowbreak{}max = 4.46 × 10⁻⁴, providing a 224× safety margin.

\textbf{Why K₇ satisfies Joyce's criterion}: The topological bound κ\_\allowbreak{}T = 1/61 constrains ‖δφ‖, placing the manifold within Joyce's perturbative regime where a torsion-free solution exists.

\subsection{11.3 Independent Numerical Validation (PINN)}\label{independent-numerical-validation-pinn}

A companion numerical program constructs explicit G₂ metrics on K₇ via physics-informed neural networks (PINNs). The three-chart atlas (neck + two Calabi-Yau bulk regions) uses approximately 10⁶ trainable parameters in float64 precision.

\textbf{Initial validation} (Phase 2):

\begin{longtable}[]{@{}lll@{}}
\toprule\noalign{}
Metric & Value & Significance \\
\midrule\noalign{}
\endhead
\bottomrule\noalign{}
\endlastfoot
‖T‖\_\allowbreak{}max & 4.46 x 10⁻⁴ & 224x below Joyce epsilon₀ \\
‖T‖\_\allowbreak{}mean & 9.8 x 10⁻⁵ & T --\textgreater{} 0 confirmed \\
det(g) error & \textless{} 10⁻⁶ & Confirms 65/32 \\
\end{longtable}

\textbf{G₂ holonomy training} (Phase 3, 13 versions, v2-v13):

Over successive training protocol refinements, the holonomy quality has improved:

\begin{longtable}[]{@{}llll@{}}
\toprule\noalign{}
Metric & Initial (v5) & Current best (v11) & Improvement \\
\midrule\noalign{}
\endhead
\bottomrule\noalign{}
\endlastfoot
g2\_\allowbreak{}self (honest holonomy) & 3.86 & 3.25 & -16\% \\
V₇ projection score & 0.51 & 0.014 & -97\% \\
det(g) at neck & 4.69 & 2.031 & locked at target \\
phi drift & 13.4\% & 0\% & controlled \\
\end{longtable}

The g2\_\allowbreak{}score measures the normalized projection of Riemann curvature onto the complement of g₂ in so(7). A score of 0 corresponds to exact G₂ holonomy; the flat metric scores approximately 3.5. The V₇ projection score measures the fraction of curvature outside the g₂ subalgebra (using the correct kernel-based g₂ decomposition, see Section 6.2).

A critical bug in the g₂ basis construction was discovered and corrected between versions 9 and 10: the Fano-plane heuristic does not produce correct g₂ generators. The correct g₂ subalgebra is the kernel of the Lie derivative map (Section 6.2). This correction led to significant improvement in all subsequent versions.

\textbf{Robust statistical validation}: The det(g) = 65/32 normalization target passes 8/8 independent consistency tests (permutation, bootstrap, Bayesian posterior 76.3\%, joint constraint p \textless{} 6 x 10⁻⁶). Per main §2.3 and §10.3 below, det(g) = 65/32 is imposed as a normalization target, not claimed as a prediction.

Full details of the PINN architecture, training protocol, and version-by-version results are presented in a companion paper.

\subsection{11.4 Lean 4 Formalization}\label{lean-4-formalization}

\textbf{Scope of verification}: The Lean formalization (146 files, 15 classified axioms of which 4 external data packages, zero \texttt{sorry}) verifies: 1. Arithmetic identities and algebraic relations between the framework's constants 2. Numerical bounds (e.g., torsion threshold) 3. G₂ differential geometry: exterior algebra Λ*(ℝ⁷), Hodge star, ψ = ⋆φ (axiom-free \texttt{Geometry} module) 4. Spectral gap bounds and topological ratios (\texttt{Spectral.\allowbreak{}Physical\allowbreak{}Spectral\allowbreak{}Gap}, 28 theorems, zero axioms). Analytical value: λ₁ = 6π²/475 = 0.12467 (see §7.3) 5. Selberg bridge: trace formula connecting S\_\allowbreak{}w(T) to spectral gap (\texttt{Spectral.\allowbreak{}Selberg\allowbreak{}Bridge}) 6. Selection principle κ = π²/14 (\texttt{Spectral.\allowbreak{}Selection\allowbreak{}Principle})

\textbf{The 4 external data-package axioms} (\texttt{K7\_\allowbreak{}analysis\_\allowbreak{}data}, \texttt{K7\_\allowbreak{}spectral\_\allowbreak{}data}, \texttt{literature\_\allowbreak{}package}, \texttt{KK\_\allowbreak{}YM\_\allowbreak{}EFT}; see main §8.1) \textbf{are used transversally throughout the formalization}: they carry the numerical and literature content (torsion bounds, the NK certificate parameter h, spectral data) that cannot be verified by \texttt{decide} alone due to real-arithmetic content. The v3.4.29 release audit classifies 15 axioms in total across an A--F taxonomy (the remaining 11 at classes A--D + F: definitional / standard / geometric / literature / numerical); main §8.1 states the reconciliation between the two counts. All other results use zero axioms.

It does \textbf{not} formalize: - Existence of K₇ as a smooth G₂ manifold - Physical interpretation of topological invariants - Uniqueness of the TCS construction

\begin{verbatim}
-- GIFT.Foundations.AnalyticalMetric

def phi0_indices : List (Fin 7 × Fin 7 × Fin 7) :=
  [(0,1,2), (0,3,4), (0,5,6), (1,3,5), (1,4,6), (2,3,6), (2,4,5)]

def phi0_signs : List Int := [1, 1, 1, 1, -1, -1, -1]

def scale_factor_power_14 : Rat := 65 / 32

theorem torsion_satisfies_joyce :
  torsion_norm_constant_form < joyce_threshold_num := by native_decide

theorem det_g_equals_target :
  scale_factor_power_14 = det_g_target := rfl
\end{verbatim}

\textbf{Status}: VERIFIED

\subsection{11.5 The Derivation Chain}\label{the-derivation-chain}

The logical structure from algebra to predictions:

\begin{verbatim}
Octonions (O)
     │
     ▼
G₂ = Aut(O), dim = 14
     │
     ▼
Standard form φ₀ (Harvey-Lawson 1982)
     │
     ▼
Scaling c = (65/32)^{1/14}    ← framework constraint
     │
     ▼
Metric g = c² × I₇
     │
     ▼
det(g) = 65/32               ← exact by imposed normalization, not fitted to experiment (§10.3)
     │
     ▼
sin²θ_W = 3/13, Q = 2/3, ...  ← Predictions
\end{verbatim}

\begin{center}\rule{\linewidth}{0.5pt}\end{center}

\section{12. Analytical G₂ Metric Details}\label{analytical-gux2082-metric-details}

\subsection{12.1 The Standard Form φ₀}\label{the-standard-form-ux3c6ux2080}

The associative 3-form preserved by G₂ ⊂ SO(7), introduced by Harvey and Lawson (1982) in their foundational work on calibrated geometries:

\[\varphi_0 = \sum_{(i,j,k) \in \mathcal{I}} \sigma_{ijk} \, e^{ijk}\]

where: - \ensuremath{\mathcal{I}} = \{(0,1,2), (0,3,4), (0,5,6), (1,3,5), (1,4,6), (2,3,6), (2,4,5)\} - σ = (+1, +1, +1, +1, -1, -1, -1)

\subsection{12.2 Linear Index Representation}\label{linear-index-representation}

In the C(7,3) = 35 basis:

\begin{longtable}[]{@{}llllll@{}}
\toprule\noalign{}
Index & Triple & Sign & Index & Triple & Sign \\
\midrule\noalign{}
\endhead
\bottomrule\noalign{}
\endlastfoot
0 & (0,1,2) & +1 & 23 & (1,4,6) & -1 \\
9 & (0,3,4) & +1 & 27 & (2,3,6) & -1 \\
14 & (0,5,6) & +1 & 28 & (2,4,5) & -1 \\
20 & (1,3,5) & +1 & -- & -- & -- \\
\end{longtable}

All other 28 components are exactly 0.

\subsection{12.3 Metric Derivation}\label{metric-derivation}

From φ₀, the metric is computed via: \[g_{ij} = \frac{1}{6} \sum_{k,l} \varphi_{ikl} \varphi_{jkl}\]

For standard φ₀: g = I₇ (identity), det(g) = 1.

Scaling φ → c·φ gives g → c²·g, hence det(g) → c¹⁴·det(g).

Setting c¹⁴ = 65/32 yields the framework metric.

\subsection{12.4 Comparison: Fano Plane vs G₂ Form}\label{comparison-fano-plane-vs-gux2082-form}

\begin{longtable}[]{@{}
  >{\raggedright\arraybackslash}p{(\columnwidth - 4\tabcolsep) * \real{0.3929}}
  >{\raggedright\arraybackslash}p{(\columnwidth - 4\tabcolsep) * \real{0.3929}}
  >{\raggedright\arraybackslash}p{(\columnwidth - 4\tabcolsep) * \real{0.2143}}@{}}
\toprule\noalign{}
\begin{minipage}[b]{\linewidth}\raggedright
Structure
\end{minipage} & \begin{minipage}[b]{\linewidth}\raggedright
7 Triples
\end{minipage} & \begin{minipage}[b]{\linewidth}\raggedright
Role
\end{minipage} \\
\midrule\noalign{}
\endhead
\bottomrule\noalign{}
\endlastfoot
\textbf{Fano lines} & (0,1,3), (1,2,4), (2,3,5), (3,4,6), (4,5,0), (5,6,1), (6,0,2) & G₂ cross-product ε\_\allowbreak{}\{ijk\} \\
\textbf{G₂ form} & (0,1,2), (0,3,4), (0,5,6), (1,3,5), (1,4,6), (2,3,6), (2,4,5) & Associative 3-form \\
\end{longtable}

Both have 7 terms but different index patterns. The Fano plane defines the octonion multiplication (cross-product), while the G₂ form is the associative calibration.

\subsection{12.5 Verification Summary}\label{verification-summary}

\begin{longtable}[]{@{}
  >{\raggedright\arraybackslash}p{(\columnwidth - 4\tabcolsep) * \real{0.2963}}
  >{\raggedright\arraybackslash}p{(\columnwidth - 4\tabcolsep) * \real{0.2963}}
  >{\raggedright\arraybackslash}p{(\columnwidth - 4\tabcolsep) * \real{0.4074}}@{}}
\toprule\noalign{}
\begin{minipage}[b]{\linewidth}\raggedright
Method
\end{minipage} & \begin{minipage}[b]{\linewidth}\raggedright
Result
\end{minipage} & \begin{minipage}[b]{\linewidth}\raggedright
Reference
\end{minipage} \\
\midrule\noalign{}
\endhead
\bottomrule\noalign{}
\endlastfoot
Algebraic & φ = (65/32)\^{}\{1/14\} × φ₀ & This section \\
Lean 4 & \texttt{det\_\allowbreak{}g\_\allowbreak{}equals\_\allowbreak{}target\ :\ rfl} & AnalyticalMetric.lean \\
PINN & Converges to constant form & gift\_\allowbreak{}core/nn/ \\
Joyce theorem & ‖T‖ \textless{} 0.1 → exists metric (224× margin) & {[}Joyce 2000{]} \\
\end{longtable}

Cross-verification between analytical and numerical methods is consistent with the conditional NK-certified solution within its analytic setup.

\begin{center}\rule{\linewidth}{0.5pt}\end{center}

\section{Author's Related Works}\label{authors-related-works}

\begin{itemize}
\tightlist
\item
  \textbf{{[}A{]}} B. de La Fournière, ``An Explicit Approximate G₂ Metric on a Compact 7-Manifold with Certified Torsion-Free Completion,'' Zenodo \href{https://doi.org/10.5281/zenodo.19892350}{10.5281/zenodo.19892350} (2026).
\item
  \textbf{{[}B{]}} B. de La Fournière, ``Spectral Geometry of the G₂-GIFT Manifold: Betti Numbers, KK Spectrum, and Spectral Invariants,'' Zenodo \href{https://doi.org/10.5281/zenodo.19893371}{10.5281/zenodo.19893371} (2026).
\item
  \textbf{{[}D{]}} B. de La Fournière, ``An Explicit Closed-Form G₂ Ansatz on a K3-Coassociative Neck with Hyperkähler Rotation and Picard-Lefschetz Wirtinger Certificate,'' Zenodo \href{https://doi.org/10.5281/zenodo.20039066}{10.5281/zenodo.20039066} (2026).
\end{itemize}

\begin{center}\rule{\linewidth}{0.5pt}\end{center}

\emph{The K₇ Framework (formerly GIFT): Supplement S1} \emph{Mathematical Foundations: E₈ + G₂ + K₇}


\bigskip
% ============================================
% GIFT v3.5 — Shared Bibliography
% ============================================
% Usage: % ============================================
% GIFT v3.5 — Shared Bibliography
% ============================================
% Usage: % ============================================
% GIFT v3.5 — Shared Bibliography
% ============================================
% Usage: \input{gift_3.5_bib} where the references section should appear.
% Sequential numbering [1]..[50]. Round-1 + round-2 IA-review updates applied
% (2026-07-07): [13] Kovalev 2005, [26] T2K/NOvA Nature 646,
% [27] NuFit 6.0 + 6.1, [45] Garbagnati-Sarti 2007, [46] LEPSUSYWG,
% [47] Fermilab DUNE 2031, [48] Chen-Hong (re-added, cited in S1), [49] Planck PR4, [50] CODATA.

\begin{thebibliography}{99}

\bibitem{1} Particle Data Group, ``Review of Particle Physics,'' \textit{Phys.\ Rev.\ D} \textbf{110}, 030001 (2024).

\bibitem{2} S.~Weinberg, ``Implications of dynamical symmetry breaking,'' \textit{Phys.\ Rev.\ D} \textbf{13}, 974 (1976).

\bibitem{3} Planck Collaboration, ``Planck 2018 results.\ VI.\ Cosmological parameters,'' \textit{Astron.\ Astrophys.} \textbf{641}, A6 (2020).

\bibitem{4} A.~G.~Riess et al., ``A comprehensive measurement of the local value of the Hubble constant,'' \textit{ApJL} \textbf{934}, L7 (2022).

\bibitem{5} C.~D.~Froggatt, H.~B.~Nielsen, ``Hierarchy of quark masses, Cabibbo angles and CP violation,'' \textit{Nucl.\ Phys.\ B} \textbf{147}, 277 (1979).

\bibitem{6} Y.~Koide, ``Fermion-boson two-body model of quarks and leptons and Cabibbo mixing,'' \textit{Lett.\ Nuovo Cim.} \textbf{34}, 201 (1982).

\bibitem{7} C.~Furey, \textit{Standard Model Physics from an Algebra?}, PhD thesis, University of Waterloo (2015).

\bibitem{8} N.~Furey, M.~J.~Hughes, ``One generation of standard model Weyl representations as a single copy of $\mathbb{R} \otimes \mathbb{C} \otimes \mathbb{H} \otimes \mathbb{O}$,'' \textit{Phys.\ Lett.\ B} \textbf{831}, 137186 (2022).

\bibitem{9} R.~Wilson, ``$E_8$ and Standard Model plus gravity,'' arXiv:2404.18938 [hep-th] (2024).

\bibitem{10} T.~P.~Singh et al., ``An $E_8 \otimes E_8$ unification of the Standard Model with pre-gravitation,'' arXiv:2206.06911v3 (2024).

\bibitem{11} B.~S.~Acharya, S.~Gukov, ``M-theory and singularities of exceptional holonomy manifolds,'' \textit{Phys.\ Rep.} \textbf{392}, 121 (2004).

\bibitem{12} L.~Foscolo et al., ``Complete non-compact $G_2$-manifolds from asymptotically conical Calabi--Yau 3-folds,'' \textit{Duke Math.\ J.} \textbf{170}, 3 (2021).

\bibitem{13} A.~Kovalev, ``Coassociative K3 fibrations of compact $G_2$-manifolds,'' arXiv:math/0511150 (2005).

\bibitem{14} M.~Haskins, J.~Nordstr\"om, ``Extra-twisted connected sum $G_2$-manifolds,'' arXiv:1809.09083 (2022); \textit{Ann.\ Glob.\ Anal.\ Geom.} \textbf{64}, art.\ 2 (2023).

\bibitem{15} G.~Bera, ``Associative submanifolds in twisted connected sum $G_2$-manifolds,'' arXiv:2209.00156 (2022).

\bibitem{16} D.~Crowley, S.~Goette, J.~Nordstr\"om, ``An analytic invariant of $G_2$ manifolds,'' \textit{Invent.\ Math.} \textbf{239}(3), 865--907 (2025).

\bibitem{17} T.~Dray, C.~A.~Manogue, \textit{The Geometry of the Octonions}, World Scientific (2015).

\bibitem{18} J.~F.~Adams, \textit{Lectures on Exceptional Lie Groups}, University of Chicago Press (1996).

\bibitem{19} D.~J.~Gross et al., ``Heterotic string theory,'' \textit{Nucl.\ Phys.\ B} \textbf{256}, 253 (1985).

\bibitem{20} D.~D.~Joyce, \textit{Compact Manifolds with Special Holonomy}, Oxford University Press (2000).

\bibitem{21} A.~Kovalev, ``Twisted connected sums and special Riemannian holonomy,'' \textit{J.\ Reine Angew.\ Math.} \textbf{565}, 125 (2003).

\bibitem{22} A.~Corti et al., ``$G_2$-manifolds and associative submanifolds via semi-Fano 3-folds,'' \textit{Duke Math.\ J.} \textbf{164}, 1971 (2015).

\bibitem{23} R.~Harvey, H.~B.~Lawson, ``Calibrated geometries,'' \textit{Acta Math.} \textbf{148}, 47 (1982).

\bibitem{24} B.~S.~Acharya, ``M-theory, $G_2$-manifolds and four-dimensional physics,'' \textit{Class.\ Quant.\ Grav.} \textbf{19}, 5619 (2002).

\bibitem{25} B.~S.~Acharya, E.~Witten, ``Chiral fermions from manifolds of $G_2$ holonomy,'' arXiv:hep-th/0109152 (2001).

\bibitem{26} T2K and NOvA Collaborations, ``Joint analysis of long-baseline neutrino oscillations,'' \textit{Nature} \textbf{646}, 818--824 (2025), \href{https://doi.org/10.1038/s41586-025-09599-3}{doi:10.1038/s41586-025-09599-3}.

\bibitem{27} I.~Esteban et al.\ (NuFIT), ``NuFit-6.0: three-flavour global analysis of neutrino oscillation experiments,'' \textit{JHEP} \textbf{12} (2024) 216, arXiv:2410.05380; NuFIT 6.1 web release, \url{https://www.nu-fit.org} (2025).

\bibitem{28} L.~de Moura, S.~Ullrich, ``The Lean 4 theorem prover and programming language,'' in \textit{CADE 28}, p.\ 625 (2021).

\bibitem{29} mathlib Community, \textit{mathlib4}, \url{https://github.com/leanprover-community/mathlib4}.

\bibitem{30} DUNE Collaboration, ``Long-baseline neutrino facility (LBNF) and DUNE conceptual design report,'' FERMILAB-TM-2696 (2020).

\bibitem{31} DUNE Collaboration, ``Deep Underground Neutrino Experiment (DUNE) far detector technical design report,'' arXiv:2103.04797 (2021).

\bibitem{32} E.~Witten, ``Strong coupling expansion of Calabi-Yau compactification,'' \textit{Nucl.\ Phys.\ B} \textbf{471}, 135 (1996).

\bibitem{33} B.~S.~Acharya et al., ``M-theory solution to the hierarchy problem,'' \textit{Phys.\ Rev.\ D} \textbf{76}, 126010 (2007).

\bibitem{34} M.~Atiyah, E.~Witten, ``M-theory dynamics on a manifold of $G_2$-holonomy,'' \textit{Adv.\ Theor.\ Math.\ Phys.} \textbf{6}, 1 (2002).

\bibitem{35} G.~Kane, \textit{String Theory and the Real World} (2017).

\bibitem{36} J.~Distler, S.~Garibaldi, ``There is no `Theory of Everything' inside $E_8$,'' \textit{Commun.\ Math.\ Phys.} \textbf{298}, 419 (2010).

\bibitem{37} J.~C.~Baez, ``Octonions and the Standard Model,'' \url{https://math.ucr.edu/home/baez/standard/} (2020--2025).

\bibitem{38} Springer Nature, ``Artificial intelligence (AI) policy,'' \url{https://www.springernature.com/gp/policies} (2024).

\bibitem{39} J.~A.~Wheeler, ``Information, physics, quantum: the search for links,'' in \textit{Complexity, Entropy, and the Physics of Information}, W.~H.~Zurek (ed.), Addison-Wesley (1990), pp.~3--28.

\bibitem{40} J.~Worrall, ``Structural realism: the best of both worlds?,'' \textit{Dialectica} \textbf{43}, 99--124 (1989).

\bibitem{41} J.~Ladyman, D.~Ross, \textit{Every Thing Must Go: Metaphysics Naturalized}, Oxford University Press (2007).

\bibitem{42} R.~Pin\v{c}\'ak, A.~Pigazzini, M.~Pudl\'ak, E.~Barto\v{s}, ``Geometric origin of a stable black hole remnant from torsion in $G_2$-manifold geometry,'' \textit{Gen.\ Rel.\ Grav.} \textbf{58}, 29 (2026), \href{https://doi.org/10.1007/s10714-026-03528-z}{doi:10.1007/s10714-026-03528-z}.

\bibitem{43} D.~Joyce, S.~Karigiannis, ``A new construction of compact torsion-free $G_2$-manifolds by gluing families of Eguchi--Hanson spaces,'' \textit{J.\ Differential Geom.} (2021), arXiv:1707.09325 (2017).

\bibitem{44} S.~Mukai, ``Finite groups of automorphisms of K3 surfaces and the Mathieu group,'' \textit{Invent.\ Math.} \textbf{94}, 183--221 (1988).

\bibitem{45} A.~Garbagnati, A.~Sarti, ``Symplectic automorphisms of prime order on K3 surfaces,'' \textit{J.\ Algebra} \textbf{318}, 323--350 (2007), arXiv:math/0603742.

\bibitem{46} LEP SUSY Working Group (ALEPH, DELPHI, L3, OPAL), ``Combined LEP chargino results, up to 208 GeV for large $m_0$,'' LEPSUSYWG/01-03.1 (2001), \url{https://lepsusy.web.cern.ch/lepsusy/www/inos_moriond01/charginos_pub.html}.

\bibitem{47} Fermilab, ``DUNE first-beam target of 2031: LBNF/DUNE-US milestone communication'' (2026-05-07), \url{https://news.fnal.gov/}.

\bibitem{48} J.~Chen, H.~Hong, ``Intermediate curvature and splitting theorem,'' arXiv:2604.26529 (2026).

\bibitem{49} M.~Tristram et al., ``Cosmological parameters derived from the final Planck data release (PR4),'' \textit{Astron.\ Astrophys.} \textbf{682}, A37 (2024).

\bibitem{50} E.~Tiesinga, P.~J.~Mohr, D.~B.~Newell, B.~N.~Taylor, ``CODATA recommended values of the fundamental physical constants: 2022,'' \textit{Rev.\ Mod.\ Phys.} \textbf{97}, 025002 (2025), \url{https://physics.nist.gov/cuu/Constants}.

\bibitem{51} V.~V.~Nikulin, ``Integer symmetric bilinear forms and some of their geometric applications,'' \textit{Izv.\ Akad.\ Nauk SSSR Ser.\ Mat.} \textbf{43} (1979); English transl.\ \textit{Math.\ USSR Izv.} \textbf{14} (1980).

\end{thebibliography}
 where the references section should appear.
% Sequential numbering [1]..[50]. Round-1 + round-2 IA-review updates applied
% (2026-07-07): [13] Kovalev 2005, [26] T2K/NOvA Nature 646,
% [27] NuFit 6.0 + 6.1, [45] Garbagnati-Sarti 2007, [46] LEPSUSYWG,
% [47] Fermilab DUNE 2031, [48] Chen-Hong (re-added, cited in S1), [49] Planck PR4, [50] CODATA.

\begin{thebibliography}{99}

\bibitem{1} Particle Data Group, ``Review of Particle Physics,'' \textit{Phys.\ Rev.\ D} \textbf{110}, 030001 (2024).

\bibitem{2} S.~Weinberg, ``Implications of dynamical symmetry breaking,'' \textit{Phys.\ Rev.\ D} \textbf{13}, 974 (1976).

\bibitem{3} Planck Collaboration, ``Planck 2018 results.\ VI.\ Cosmological parameters,'' \textit{Astron.\ Astrophys.} \textbf{641}, A6 (2020).

\bibitem{4} A.~G.~Riess et al., ``A comprehensive measurement of the local value of the Hubble constant,'' \textit{ApJL} \textbf{934}, L7 (2022).

\bibitem{5} C.~D.~Froggatt, H.~B.~Nielsen, ``Hierarchy of quark masses, Cabibbo angles and CP violation,'' \textit{Nucl.\ Phys.\ B} \textbf{147}, 277 (1979).

\bibitem{6} Y.~Koide, ``Fermion-boson two-body model of quarks and leptons and Cabibbo mixing,'' \textit{Lett.\ Nuovo Cim.} \textbf{34}, 201 (1982).

\bibitem{7} C.~Furey, \textit{Standard Model Physics from an Algebra?}, PhD thesis, University of Waterloo (2015).

\bibitem{8} N.~Furey, M.~J.~Hughes, ``One generation of standard model Weyl representations as a single copy of $\mathbb{R} \otimes \mathbb{C} \otimes \mathbb{H} \otimes \mathbb{O}$,'' \textit{Phys.\ Lett.\ B} \textbf{831}, 137186 (2022).

\bibitem{9} R.~Wilson, ``$E_8$ and Standard Model plus gravity,'' arXiv:2404.18938 [hep-th] (2024).

\bibitem{10} T.~P.~Singh et al., ``An $E_8 \otimes E_8$ unification of the Standard Model with pre-gravitation,'' arXiv:2206.06911v3 (2024).

\bibitem{11} B.~S.~Acharya, S.~Gukov, ``M-theory and singularities of exceptional holonomy manifolds,'' \textit{Phys.\ Rep.} \textbf{392}, 121 (2004).

\bibitem{12} L.~Foscolo et al., ``Complete non-compact $G_2$-manifolds from asymptotically conical Calabi--Yau 3-folds,'' \textit{Duke Math.\ J.} \textbf{170}, 3 (2021).

\bibitem{13} A.~Kovalev, ``Coassociative K3 fibrations of compact $G_2$-manifolds,'' arXiv:math/0511150 (2005).

\bibitem{14} M.~Haskins, J.~Nordstr\"om, ``Extra-twisted connected sum $G_2$-manifolds,'' arXiv:1809.09083 (2022); \textit{Ann.\ Glob.\ Anal.\ Geom.} \textbf{64}, art.\ 2 (2023).

\bibitem{15} G.~Bera, ``Associative submanifolds in twisted connected sum $G_2$-manifolds,'' arXiv:2209.00156 (2022).

\bibitem{16} D.~Crowley, S.~Goette, J.~Nordstr\"om, ``An analytic invariant of $G_2$ manifolds,'' \textit{Invent.\ Math.} \textbf{239}(3), 865--907 (2025).

\bibitem{17} T.~Dray, C.~A.~Manogue, \textit{The Geometry of the Octonions}, World Scientific (2015).

\bibitem{18} J.~F.~Adams, \textit{Lectures on Exceptional Lie Groups}, University of Chicago Press (1996).

\bibitem{19} D.~J.~Gross et al., ``Heterotic string theory,'' \textit{Nucl.\ Phys.\ B} \textbf{256}, 253 (1985).

\bibitem{20} D.~D.~Joyce, \textit{Compact Manifolds with Special Holonomy}, Oxford University Press (2000).

\bibitem{21} A.~Kovalev, ``Twisted connected sums and special Riemannian holonomy,'' \textit{J.\ Reine Angew.\ Math.} \textbf{565}, 125 (2003).

\bibitem{22} A.~Corti et al., ``$G_2$-manifolds and associative submanifolds via semi-Fano 3-folds,'' \textit{Duke Math.\ J.} \textbf{164}, 1971 (2015).

\bibitem{23} R.~Harvey, H.~B.~Lawson, ``Calibrated geometries,'' \textit{Acta Math.} \textbf{148}, 47 (1982).

\bibitem{24} B.~S.~Acharya, ``M-theory, $G_2$-manifolds and four-dimensional physics,'' \textit{Class.\ Quant.\ Grav.} \textbf{19}, 5619 (2002).

\bibitem{25} B.~S.~Acharya, E.~Witten, ``Chiral fermions from manifolds of $G_2$ holonomy,'' arXiv:hep-th/0109152 (2001).

\bibitem{26} T2K and NOvA Collaborations, ``Joint analysis of long-baseline neutrino oscillations,'' \textit{Nature} \textbf{646}, 818--824 (2025), \href{https://doi.org/10.1038/s41586-025-09599-3}{doi:10.1038/s41586-025-09599-3}.

\bibitem{27} I.~Esteban et al.\ (NuFIT), ``NuFit-6.0: three-flavour global analysis of neutrino oscillation experiments,'' \textit{JHEP} \textbf{12} (2024) 216, arXiv:2410.05380; NuFIT 6.1 web release, \url{https://www.nu-fit.org} (2025).

\bibitem{28} L.~de Moura, S.~Ullrich, ``The Lean 4 theorem prover and programming language,'' in \textit{CADE 28}, p.\ 625 (2021).

\bibitem{29} mathlib Community, \textit{mathlib4}, \url{https://github.com/leanprover-community/mathlib4}.

\bibitem{30} DUNE Collaboration, ``Long-baseline neutrino facility (LBNF) and DUNE conceptual design report,'' FERMILAB-TM-2696 (2020).

\bibitem{31} DUNE Collaboration, ``Deep Underground Neutrino Experiment (DUNE) far detector technical design report,'' arXiv:2103.04797 (2021).

\bibitem{32} E.~Witten, ``Strong coupling expansion of Calabi-Yau compactification,'' \textit{Nucl.\ Phys.\ B} \textbf{471}, 135 (1996).

\bibitem{33} B.~S.~Acharya et al., ``M-theory solution to the hierarchy problem,'' \textit{Phys.\ Rev.\ D} \textbf{76}, 126010 (2007).

\bibitem{34} M.~Atiyah, E.~Witten, ``M-theory dynamics on a manifold of $G_2$-holonomy,'' \textit{Adv.\ Theor.\ Math.\ Phys.} \textbf{6}, 1 (2002).

\bibitem{35} G.~Kane, \textit{String Theory and the Real World} (2017).

\bibitem{36} J.~Distler, S.~Garibaldi, ``There is no `Theory of Everything' inside $E_8$,'' \textit{Commun.\ Math.\ Phys.} \textbf{298}, 419 (2010).

\bibitem{37} J.~C.~Baez, ``Octonions and the Standard Model,'' \url{https://math.ucr.edu/home/baez/standard/} (2020--2025).

\bibitem{38} Springer Nature, ``Artificial intelligence (AI) policy,'' \url{https://www.springernature.com/gp/policies} (2024).

\bibitem{39} J.~A.~Wheeler, ``Information, physics, quantum: the search for links,'' in \textit{Complexity, Entropy, and the Physics of Information}, W.~H.~Zurek (ed.), Addison-Wesley (1990), pp.~3--28.

\bibitem{40} J.~Worrall, ``Structural realism: the best of both worlds?,'' \textit{Dialectica} \textbf{43}, 99--124 (1989).

\bibitem{41} J.~Ladyman, D.~Ross, \textit{Every Thing Must Go: Metaphysics Naturalized}, Oxford University Press (2007).

\bibitem{42} R.~Pin\v{c}\'ak, A.~Pigazzini, M.~Pudl\'ak, E.~Barto\v{s}, ``Geometric origin of a stable black hole remnant from torsion in $G_2$-manifold geometry,'' \textit{Gen.\ Rel.\ Grav.} \textbf{58}, 29 (2026), \href{https://doi.org/10.1007/s10714-026-03528-z}{doi:10.1007/s10714-026-03528-z}.

\bibitem{43} D.~Joyce, S.~Karigiannis, ``A new construction of compact torsion-free $G_2$-manifolds by gluing families of Eguchi--Hanson spaces,'' \textit{J.\ Differential Geom.} (2021), arXiv:1707.09325 (2017).

\bibitem{44} S.~Mukai, ``Finite groups of automorphisms of K3 surfaces and the Mathieu group,'' \textit{Invent.\ Math.} \textbf{94}, 183--221 (1988).

\bibitem{45} A.~Garbagnati, A.~Sarti, ``Symplectic automorphisms of prime order on K3 surfaces,'' \textit{J.\ Algebra} \textbf{318}, 323--350 (2007), arXiv:math/0603742.

\bibitem{46} LEP SUSY Working Group (ALEPH, DELPHI, L3, OPAL), ``Combined LEP chargino results, up to 208 GeV for large $m_0$,'' LEPSUSYWG/01-03.1 (2001), \url{https://lepsusy.web.cern.ch/lepsusy/www/inos_moriond01/charginos_pub.html}.

\bibitem{47} Fermilab, ``DUNE first-beam target of 2031: LBNF/DUNE-US milestone communication'' (2026-05-07), \url{https://news.fnal.gov/}.

\bibitem{48} J.~Chen, H.~Hong, ``Intermediate curvature and splitting theorem,'' arXiv:2604.26529 (2026).

\bibitem{49} M.~Tristram et al., ``Cosmological parameters derived from the final Planck data release (PR4),'' \textit{Astron.\ Astrophys.} \textbf{682}, A37 (2024).

\bibitem{50} E.~Tiesinga, P.~J.~Mohr, D.~B.~Newell, B.~N.~Taylor, ``CODATA recommended values of the fundamental physical constants: 2022,'' \textit{Rev.\ Mod.\ Phys.} \textbf{97}, 025002 (2025), \url{https://physics.nist.gov/cuu/Constants}.

\bibitem{51} V.~V.~Nikulin, ``Integer symmetric bilinear forms and some of their geometric applications,'' \textit{Izv.\ Akad.\ Nauk SSSR Ser.\ Mat.} \textbf{43} (1979); English transl.\ \textit{Math.\ USSR Izv.} \textbf{14} (1980).

\end{thebibliography}
 where the references section should appear.
% Sequential numbering [1]..[50]. Round-1 + round-2 IA-review updates applied
% (2026-07-07): [13] Kovalev 2005, [26] T2K/NOvA Nature 646,
% [27] NuFit 6.0 + 6.1, [45] Garbagnati-Sarti 2007, [46] LEPSUSYWG,
% [47] Fermilab DUNE 2031, [48] Chen-Hong (re-added, cited in S1), [49] Planck PR4, [50] CODATA.

\begin{thebibliography}{99}

\bibitem{1} Particle Data Group, ``Review of Particle Physics,'' \textit{Phys.\ Rev.\ D} \textbf{110}, 030001 (2024).

\bibitem{2} S.~Weinberg, ``Implications of dynamical symmetry breaking,'' \textit{Phys.\ Rev.\ D} \textbf{13}, 974 (1976).

\bibitem{3} Planck Collaboration, ``Planck 2018 results.\ VI.\ Cosmological parameters,'' \textit{Astron.\ Astrophys.} \textbf{641}, A6 (2020).

\bibitem{4} A.~G.~Riess et al., ``A comprehensive measurement of the local value of the Hubble constant,'' \textit{ApJL} \textbf{934}, L7 (2022).

\bibitem{5} C.~D.~Froggatt, H.~B.~Nielsen, ``Hierarchy of quark masses, Cabibbo angles and CP violation,'' \textit{Nucl.\ Phys.\ B} \textbf{147}, 277 (1979).

\bibitem{6} Y.~Koide, ``Fermion-boson two-body model of quarks and leptons and Cabibbo mixing,'' \textit{Lett.\ Nuovo Cim.} \textbf{34}, 201 (1982).

\bibitem{7} C.~Furey, \textit{Standard Model Physics from an Algebra?}, PhD thesis, University of Waterloo (2015).

\bibitem{8} N.~Furey, M.~J.~Hughes, ``One generation of standard model Weyl representations as a single copy of $\mathbb{R} \otimes \mathbb{C} \otimes \mathbb{H} \otimes \mathbb{O}$,'' \textit{Phys.\ Lett.\ B} \textbf{831}, 137186 (2022).

\bibitem{9} R.~Wilson, ``$E_8$ and Standard Model plus gravity,'' arXiv:2404.18938 [hep-th] (2024).

\bibitem{10} T.~P.~Singh et al., ``An $E_8 \otimes E_8$ unification of the Standard Model with pre-gravitation,'' arXiv:2206.06911v3 (2024).

\bibitem{11} B.~S.~Acharya, S.~Gukov, ``M-theory and singularities of exceptional holonomy manifolds,'' \textit{Phys.\ Rep.} \textbf{392}, 121 (2004).

\bibitem{12} L.~Foscolo et al., ``Complete non-compact $G_2$-manifolds from asymptotically conical Calabi--Yau 3-folds,'' \textit{Duke Math.\ J.} \textbf{170}, 3 (2021).

\bibitem{13} A.~Kovalev, ``Coassociative K3 fibrations of compact $G_2$-manifolds,'' arXiv:math/0511150 (2005).

\bibitem{14} M.~Haskins, J.~Nordstr\"om, ``Extra-twisted connected sum $G_2$-manifolds,'' arXiv:1809.09083 (2022); \textit{Ann.\ Glob.\ Anal.\ Geom.} \textbf{64}, art.\ 2 (2023).

\bibitem{15} G.~Bera, ``Associative submanifolds in twisted connected sum $G_2$-manifolds,'' arXiv:2209.00156 (2022).

\bibitem{16} D.~Crowley, S.~Goette, J.~Nordstr\"om, ``An analytic invariant of $G_2$ manifolds,'' \textit{Invent.\ Math.} \textbf{239}(3), 865--907 (2025).

\bibitem{17} T.~Dray, C.~A.~Manogue, \textit{The Geometry of the Octonions}, World Scientific (2015).

\bibitem{18} J.~F.~Adams, \textit{Lectures on Exceptional Lie Groups}, University of Chicago Press (1996).

\bibitem{19} D.~J.~Gross et al., ``Heterotic string theory,'' \textit{Nucl.\ Phys.\ B} \textbf{256}, 253 (1985).

\bibitem{20} D.~D.~Joyce, \textit{Compact Manifolds with Special Holonomy}, Oxford University Press (2000).

\bibitem{21} A.~Kovalev, ``Twisted connected sums and special Riemannian holonomy,'' \textit{J.\ Reine Angew.\ Math.} \textbf{565}, 125 (2003).

\bibitem{22} A.~Corti et al., ``$G_2$-manifolds and associative submanifolds via semi-Fano 3-folds,'' \textit{Duke Math.\ J.} \textbf{164}, 1971 (2015).

\bibitem{23} R.~Harvey, H.~B.~Lawson, ``Calibrated geometries,'' \textit{Acta Math.} \textbf{148}, 47 (1982).

\bibitem{24} B.~S.~Acharya, ``M-theory, $G_2$-manifolds and four-dimensional physics,'' \textit{Class.\ Quant.\ Grav.} \textbf{19}, 5619 (2002).

\bibitem{25} B.~S.~Acharya, E.~Witten, ``Chiral fermions from manifolds of $G_2$ holonomy,'' arXiv:hep-th/0109152 (2001).

\bibitem{26} T2K and NOvA Collaborations, ``Joint analysis of long-baseline neutrino oscillations,'' \textit{Nature} \textbf{646}, 818--824 (2025), \href{https://doi.org/10.1038/s41586-025-09599-3}{doi:10.1038/s41586-025-09599-3}.

\bibitem{27} I.~Esteban et al.\ (NuFIT), ``NuFit-6.0: three-flavour global analysis of neutrino oscillation experiments,'' \textit{JHEP} \textbf{12} (2024) 216, arXiv:2410.05380; NuFIT 6.1 web release, \url{https://www.nu-fit.org} (2025).

\bibitem{28} L.~de Moura, S.~Ullrich, ``The Lean 4 theorem prover and programming language,'' in \textit{CADE 28}, p.\ 625 (2021).

\bibitem{29} mathlib Community, \textit{mathlib4}, \url{https://github.com/leanprover-community/mathlib4}.

\bibitem{30} DUNE Collaboration, ``Long-baseline neutrino facility (LBNF) and DUNE conceptual design report,'' FERMILAB-TM-2696 (2020).

\bibitem{31} DUNE Collaboration, ``Deep Underground Neutrino Experiment (DUNE) far detector technical design report,'' arXiv:2103.04797 (2021).

\bibitem{32} E.~Witten, ``Strong coupling expansion of Calabi-Yau compactification,'' \textit{Nucl.\ Phys.\ B} \textbf{471}, 135 (1996).

\bibitem{33} B.~S.~Acharya et al., ``M-theory solution to the hierarchy problem,'' \textit{Phys.\ Rev.\ D} \textbf{76}, 126010 (2007).

\bibitem{34} M.~Atiyah, E.~Witten, ``M-theory dynamics on a manifold of $G_2$-holonomy,'' \textit{Adv.\ Theor.\ Math.\ Phys.} \textbf{6}, 1 (2002).

\bibitem{35} G.~Kane, \textit{String Theory and the Real World} (2017).

\bibitem{36} J.~Distler, S.~Garibaldi, ``There is no `Theory of Everything' inside $E_8$,'' \textit{Commun.\ Math.\ Phys.} \textbf{298}, 419 (2010).

\bibitem{37} J.~C.~Baez, ``Octonions and the Standard Model,'' \url{https://math.ucr.edu/home/baez/standard/} (2020--2025).

\bibitem{38} Springer Nature, ``Artificial intelligence (AI) policy,'' \url{https://www.springernature.com/gp/policies} (2024).

\bibitem{39} J.~A.~Wheeler, ``Information, physics, quantum: the search for links,'' in \textit{Complexity, Entropy, and the Physics of Information}, W.~H.~Zurek (ed.), Addison-Wesley (1990), pp.~3--28.

\bibitem{40} J.~Worrall, ``Structural realism: the best of both worlds?,'' \textit{Dialectica} \textbf{43}, 99--124 (1989).

\bibitem{41} J.~Ladyman, D.~Ross, \textit{Every Thing Must Go: Metaphysics Naturalized}, Oxford University Press (2007).

\bibitem{42} R.~Pin\v{c}\'ak, A.~Pigazzini, M.~Pudl\'ak, E.~Barto\v{s}, ``Geometric origin of a stable black hole remnant from torsion in $G_2$-manifold geometry,'' \textit{Gen.\ Rel.\ Grav.} \textbf{58}, 29 (2026), \href{https://doi.org/10.1007/s10714-026-03528-z}{doi:10.1007/s10714-026-03528-z}.

\bibitem{43} D.~Joyce, S.~Karigiannis, ``A new construction of compact torsion-free $G_2$-manifolds by gluing families of Eguchi--Hanson spaces,'' \textit{J.\ Differential Geom.} (2021), arXiv:1707.09325 (2017).

\bibitem{44} S.~Mukai, ``Finite groups of automorphisms of K3 surfaces and the Mathieu group,'' \textit{Invent.\ Math.} \textbf{94}, 183--221 (1988).

\bibitem{45} A.~Garbagnati, A.~Sarti, ``Symplectic automorphisms of prime order on K3 surfaces,'' \textit{J.\ Algebra} \textbf{318}, 323--350 (2007), arXiv:math/0603742.

\bibitem{46} LEP SUSY Working Group (ALEPH, DELPHI, L3, OPAL), ``Combined LEP chargino results, up to 208 GeV for large $m_0$,'' LEPSUSYWG/01-03.1 (2001), \url{https://lepsusy.web.cern.ch/lepsusy/www/inos_moriond01/charginos_pub.html}.

\bibitem{47} Fermilab, ``DUNE first-beam target of 2031: LBNF/DUNE-US milestone communication'' (2026-05-07), \url{https://news.fnal.gov/}.

\bibitem{48} J.~Chen, H.~Hong, ``Intermediate curvature and splitting theorem,'' arXiv:2604.26529 (2026).

\bibitem{49} M.~Tristram et al., ``Cosmological parameters derived from the final Planck data release (PR4),'' \textit{Astron.\ Astrophys.} \textbf{682}, A37 (2024).

\bibitem{50} E.~Tiesinga, P.~J.~Mohr, D.~B.~Newell, B.~N.~Taylor, ``CODATA recommended values of the fundamental physical constants: 2022,'' \textit{Rev.\ Mod.\ Phys.} \textbf{97}, 025002 (2025), \url{https://physics.nist.gov/cuu/Constants}.

\bibitem{51} V.~V.~Nikulin, ``Integer symmetric bilinear forms and some of their geometric applications,'' \textit{Izv.\ Akad.\ Nauk SSSR Ser.\ Mat.} \textbf{43} (1979); English transl.\ \textit{Math.\ USSR Izv.} \textbf{14} (1980).

\end{thebibliography}


\end{document}
